\documentclass[aspectratio=169,10pt]{beamer}
% notes.tex defines \NOTESONLY and reads this file: then only the speaker notes are printed
\ifdefined\NOTESONLY
  \setbeameroption{show only notes}
\fi

\usepackage{lmodern}
\usepackage[utf8]{inputenc}
\usepackage[T1]{fontenc}
\usepackage[polish,english]{babel}
\usepackage{amsmath,amssymb}
\usepackage{booktabs}
\usepackage{colortbl}
\usepackage{xcolor}
\usepackage{fontawesome5}
\usepackage{listings}
\usepackage{etoolbox}
\usepackage{xstring}
\usepackage{tikz}
\usetikzlibrary{arrows.meta, calc, positioning, fit, decorations.pathreplacing, decorations.pathmorphing}
\usepackage{pgfplots}
% layers for the code windows drawn in TikZ: the window, then the line that runs now, then the text
\pgfdeclarelayer{codebg}\pgfdeclarelayer{codehl}\pgfsetlayers{codebg,codehl,main}
\pgfplotsset{compat=1.18}

\usetheme{Madrid}
\useinnertheme{rounded}
\useoutertheme{miniframes}
\setbeamertemplate{navigation symbols}{}
\setbeamertemplate{itemize items}[circle]
\setbeamertemplate{enumerate items}[default]
\setbeamertemplate{blocks}[rounded][shadow=true]
% one dot per idea, not per frame: \samedot before a frame gives it the dot of the frame before
% (it undoes the step of the dot counter that the next frame makes)
\newcommand{\samedot}{\addtocounter{subsectionslide}{-1}}
% the default dot step of the circle dots (0.1 cm dots)
\setbeamersize{mini frame size=.14cm,mini frame offset=-.01cm}

% ======================================================================
% Colour rules for presentations (Beamer + TikZ), from Agents/colors
% ======================================================================

% ---------- template ----------
\definecolor{DeepBlue}{HTML}{283593}      % title bar, headings, links, control-flow arrows
\definecolor{AccentOrange}{HTML}{EB5A14}  % annotations: step numbers, "look here", code highlight, conclusions
\definecolor{LightGray}{HTML}{F5F6FA}     % background of information blocks
\setbeamercolor{palette primary}{bg=DeepBlue,fg=white}
\setbeamercolor{palette secondary}{bg=DeepBlue!80!black,fg=white}
\setbeamercolor{palette tertiary}{bg=DeepBlue!90!black,fg=white}
\setbeamercolor{title}{bg=DeepBlue,fg=white}
\setbeamercolor{frametitle}{bg=DeepBlue,fg=white}
\setbeamercolor{structure}{fg=DeepBlue}
\setbeamercolor{item}{fg=AccentOrange}
\setbeamercolor{alerted text}{fg=AccentOrange}
\setbeamercolor{block title}{bg=DeepBlue,fg=white}
\setbeamercolor{block body}{bg=LightGray,fg=black}
\setbeamercolor{block title example}{bg=AccentOrange,fg=white}
\setbeamercolor{block body example}{bg=AccentOrange!10,fg=black}
\setbeamercolor{enumerate item}{fg=DeepBlue}
\setbeamerfont{enumerate item}{series=\bfseries}

% ---------- actors and categories: fills only (S1 to S6) ----------
\definecolor{sIndigo}{HTML}{5C6BC0}   % S1
\definecolor{sPink}{HTML}{EC407A}     % S2
\definecolor{sViolet}{HTML}{7E57C2}   % S3
\definecolor{sAmber}{HTML}{FFB300}    % S4: black text on it
\definecolor{sTeal}{HTML}{26A69A}     % S5
\definecolor{sBrown}{HTML}{8D6E63}    % S6
% text variants that stay readable on white
\definecolor{sPinkText}{HTML}{C2185B}
\definecolor{sAmberText}{HTML}{A06A00}
\definecolor{sTealText}{HTML}{00796B}

% ---------- actions: frames only; results: signs only ----------
\definecolor{aExec}{HTML}{43A047}     % runs now
\definecolor{aRead}{HTML}{1E88E5}     % is read
\definecolor{aWrite}{HTML}{E53935}    % is written

% ---------- neutral (Blue Grey) ----------
\definecolor{nInk}{HTML}{37474F}      % text in figures, waveforms, wires, data-flow arrows
\definecolor{nLine}{HTML}{90A4AE}     % borders, axes
\definecolor{nLight}{HTML}{CFD8DC}    % inactive lines, grids, empty places (dashed)
\definecolor{nFill}{HTML}{ECEFF1}     % fill of boxes and used cells
\definecolor{nIdle}{HTML}{E0E0E0}     % time lost (idle, busy waiting, bubbles)

% ---------- TikZ styles ----------
\tikzset{
  aexec/.style={draw=aExec, line width=1.6pt},
  aexecm/.style={aexec, label={[text=aExec, inner sep=1pt, font=\tiny]left:$\blacktriangleright$}},
  aread/.style={draw=aRead, line width=1.6pt, dash pattern=on 3pt off 2pt},
  awrite/.style={draw=aWrite, line width=2.4pt},
  hw/.style={draw=nLine, thick, fill=nFill, text=nInk, rounded corners=3pt, align=center, inner sep=4pt},
  hwmain/.style={draw=nInk, thick, fill=nInk, text=white, rounded corners=3pt, align=center, inner sep=4pt},
  cell/.style={draw=nLine, thick, fill=white, text=nInk, font=\ttfamily\scriptsize},
  cellused/.style={cell, fill=nFill},
  cellfree/.style={cell, dashed, draw=nLight},
  container/.style={draw=nLine, thick, dashed, rounded corners=5pt},
  actor/.style={fill=#1, draw=#1, text=white, rounded corners=3pt, align=center, font=\scriptsize\bfseries, inner sep=4pt},
  actorlight/.style={fill=#1!15, draw=#1, thick, text=nInk, rounded corners=3pt, align=center, inner sep=4pt},
  suspended/.style={fill=#1!30, draw=#1, dashed, line width=0.9pt, text=nInk, rounded corners=3pt, align=center, inner sep=4pt},
  ready/.style={fill=white, draw=#1, line width=0.8pt, text=nInk, rounded corners=3pt, align=center, inner sep=4pt},
  idle/.style={fill=nIdle, draw=none, text=nInk, rounded corners=3pt, align=center, inner sep=4pt},
  flow/.style={-{Stealth}, very thick, DeepBlue},
  data/.style={-{Stealth}, very thick, nInk},
  holds/.style={-{Stealth}, very thick, nInk},
  waits/.style={-{Stealth}, very thick, nLine, dashed},
  annot/.style={draw=AccentOrange, line width=1.4pt, rounded corners=3pt},
}
\tikzset{actoramber/.style={actor=sAmber, text=black}}

% ---------- signs and small helpers ----------
\newcommand{\ok}{\textcolor{aExec}{\faCheck}}
\newcommand{\nok}{\textcolor{aWrite}{\faTimes}}
\newcommand{\badge}[1]{\tikz[baseline=(b.base)]\node[circle,fill=AccentOrange,text=white,
  inner sep=1.2pt,font=\bfseries\tiny] (b) {#1};}
\newcommand{\bigbadge}[1]{\tikz[baseline=(b.base)]\node[circle,fill=AccentOrange,text=white,
  inner sep=2pt,font=\bfseries\small] (b) {#1};}
% ---------- text boxes ----------
\newcommand{\defbox}[1]{%
  \begin{beamercolorbox}[rounded=true,sep=0.1cm,wd=\linewidth]{block body example}
    \footnotesize #1
  \end{beamercolorbox}\vspace{0.06cm}}
% a large definition in the middle of the slide
\newcommand{\bigdef}[1]{%
  \begin{beamercolorbox}[rounded=true,sep=0.25cm,wd=\linewidth]{block body example}
    \centering\normalsize #1
  \end{beamercolorbox}\vspace{0.06cm}}
\newcommand{\stepbox}[2]{%
  \begin{beamercolorbox}[rounded=true,sep=0.09cm,wd=\linewidth]{block body}
    \footnotesize\textbf{#1}\\ #2
  \end{beamercolorbox}\vspace{0.06cm}}
\newcommand{\keymsg}[1]{%
  \begin{beamercolorbox}[rounded=true,sep=0.1cm,wd=\textwidth]{block body example}
    \centering\footnotesize #1
  \end{beamercolorbox}}
% a grey box without a title
\newcommand{\plainbox}[1]{%
  \begin{beamercolorbox}[rounded=true,sep=0.12cm,wd=\linewidth]{block body}
    \centering\footnotesize #1
  \end{beamercolorbox}\vspace{0.06cm}}
\setbeamercolor{block body problem}{bg=AccentOrange!10,fg=black}
\newcommand{\probbox}[2]{%
  \begin{beamercolorbox}[rounded=true,sep=0.09cm,wd=\linewidth]{block body problem}
    \footnotesize\textbf{\textcolor{AccentOrange}{\faExclamationTriangle}\ #1}\\ #2
  \end{beamercolorbox}\vspace{0.06cm}}

% ---------- code: dark window, grey dots, syntax colours ----------
\definecolor{VSDarkBg}{RGB}{25,25,25}
\definecolor{VSText}{RGB}{245,245,245}
\definecolor{VSBlue}{RGB}{80,180,255}
\definecolor{VSPurple}{RGB}{210,110,210}
\definecolor{VSString}{RGB}{255,140,80}
\definecolor{VSComment}{RGB}{80,220,80}
\setbeamercolor{vscodelisting}{bg=VSDarkBg}
\lstdefinestyle{code}{basicstyle=\color{VSText}\ttfamily\scriptsize, keywordstyle=\color{VSBlue}\bfseries,
  stringstyle=\color{VSString}, commentstyle=\color{VSComment}\itshape, emphstyle=\color{VSPurple},
  numbers=none, columns=fixed, basewidth=0.52em, keepspaces=true, showstringspaces=false, xleftmargin=6pt}
\lstdefinestyle{py}{style=code, language=Python, emph={print,chr,zip,range,len,bytes,ord,hex}}
\BeforeBeginEnvironment{lstlisting}{%
  \begin{beamercolorbox}[rounded=true,sep=2mm,wd=\linewidth]{vscodelisting}%
  \begin{tikzpicture}
    \fill[nLine] (0,0) circle (3pt); \fill[nLine] (0.35,0) circle (3pt); \fill[nLine] (0.70,0) circle (3pt);
  \end{tikzpicture}%
  \vspace{-0.2cm}}
\AfterEndEnvironment{lstlisting}{\end{beamercolorbox}}
\newcommand{\hlcode}[1]{{\setlength{\fboxsep}{1pt}\colorbox{AccentOrange}{\textcolor{white}{\bfseries #1}}}}
\newcommand{\eol}{\tikz[remember picture, overlay]\coordinate (eol) at (0,0.08);}
\newcommand{\codenote}[3][0pt]{%
  \begin{tikzpicture}[remember picture, overlay, >=Stealth]
    \node[draw=AccentOrange, very thick, fill=AccentOrange!6, rounded corners=4pt, text width=0.33\paperwidth,
      align=left, inner sep=6pt, font=\footnotesize, anchor=west] (cal)
      at ([yshift=#1]eol -| {$(current page.west)!0.6!(current page.east)$}) {\textbf{\textcolor{AccentOrange}{#2}}\\[3pt]#3};
    \draw[->, very thick, AccentOrange] (eol -| cal.west) -- ([xshift=4pt]eol);
  \end{tikzpicture}}
% program output: a light terminal line
\newcommand{\terminal}[1]{%
  \begin{tikzpicture}
    \node[draw=nLine, thick, fill=nFill, rounded corners=3pt, inner sep=5pt, font=\ttfamily\small, text=nInk]
      {\textcolor{nLine}{\$}\ #1};
  \end{tikzpicture}}

% ======================================================================
% Course macros
% ======================================================================
% task header: \taskhead{Task}{Open the template from \ghlink{Task1}.}
\newcommand{\taskhead}[2]{%
  \begin{beamercolorbox}[rounded=true,sep=0.12cm,wd=\linewidth]{block body example}
    \small\textbf{#1}\\[0.04cm]
    \footnotesize #2
  \end{beamercolorbox}\vspace{0.1cm}}
% the starter files on GitHub: the word "GitHub" is the link.
%   \ghlink{Task4} opens the folder of Task 4, \ghlink{} opens the folder of this lecture
\newcommand{\ghbase}{https://github.com/KsaweryM/Cryptography-Course/tree/main/Lecture2}
\newcommand{\ghlink}[1]{\href{\ghbase/#1}{\textcolor{DeepBlue}{\textbf{GitHub}}}}

% speaker notes (Polish): printed only in the notes PDF (make notes)
\setbeamertemplate{note page}[default]
\setbeamerfont{note page}{size=\normalsize}
\newcommand{\say}[1]{\note{\selectlanguage{polish}#1}}

% ======================================================================
% Colour mapping for this lecture (colour rules, section 6)
%   S1 sIndigo : the sender (the sensor)
%   S2 sPink   : the attacker
%   S3 sViolet : plaintext
%   S4 sAmber  : key and keystream
%   S5 sTeal   : the receiver (the server)
%   S6 sBrown  : ciphertext
%   mechanisms (generator, block cipher, XOR, RC4 itself) are neutral
% ======================================================================
\tikzset{
  byte/.style={draw=nLine, thick, fill=white, text=nInk, font=\ttfamily\scriptsize,
    minimum width=0.62cm, minimum height=0.48cm, inner sep=1pt, rounded corners=1.5pt},
  ptb/.style={byte, fill=sViolet!15, draw=sViolet},
  ksb/.style={byte, fill=sAmber!25, draw=sAmber},
  ctb/.style={byte, fill=sBrown!15, draw=sBrown},
  % the same ciphertext value always gets the same shade of S6 (lightness, not a new colour)
  ct58/.style={ctb, fill=sBrown!12},
  ct137/.style={ctb, fill=sBrown!35},
  ct201/.style={ctb, fill=sBrown!70, text=white},
  ct64/.style={ctb, fill=sBrown!100, text=white},
  rowlabel/.style={font=\scriptsize, text=nInk, anchor=east},
  sender/.style={actor=sIndigo, minimum width=2.4cm, minimum height=1.0cm, font=\small\bfseries},
  receiver/.style={actor=sTeal, minimum width=2.4cm, minimum height=1.0cm, font=\small\bfseries},
  attacker/.style={actor=sPink, minimum width=2.0cm, minimum height=0.8cm, font=\small\bfseries},
  radio/.style={-{Stealth}, very thick, nInk, decorate,
    decoration={snake, amplitude=1.2pt, segment length=7pt, post length=4pt}},
}
% a row of bytes: \byterow{x of the first cell}{y}{style}{list of values}{step between cells}
\newcommand{\byterow}[5]{%
  \foreach \v [count=\k from 0] in {#4} {
    \node[#3] at ({#1+#5*\k},#2) {\v};
  }}


% ======================================================================
% Helpers shared with Lecture 1 (same look in every lecture)
% ======================================================================
% an intercepted envelope with ciphertext: \envelope{name}{position}{label}
\newcommand{\envelope}[3]{%
  \node[draw=sBrown, thick, fill=sBrown!15, minimum width=2.2cm, minimum height=1.3cm, rounded corners=1pt] (#1) at (#2) {};
  \draw[sBrown, thick] (#1.north west) -- ([yshift=-0.45cm]#1.north) -- (#1.north east);
  \node[circle, fill=sBrown, text=white, inner sep=1.5pt, font=\scriptsize] at ([yshift=-0.45cm]#1.north) {\faLock};
  \node[font=\scriptsize\bfseries, text=nInk] at ([yshift=-0.28cm]#1.center) {#3};}
% a note with plaintext on it: \papernote{name}{position}{text}
\newcommand{\papernote}[3]{%
  \node[draw=sViolet, thick, fill=sViolet!12, rounded corners=1pt, inner sep=6pt, font=\ttfamily\small, text=nInk,
    rotate=3] (#1) at (#2) {#3};}
% code drawn in TikZ, so that single lines can be marked:
%   \codebox{number of lines}{width}  then  \cl{line}{text}, \clmark{line} (runs now), \clnote{line}{text}
\newcommand{\codepitch}{0.36}
\newcommand{\codebox}[2]{%
  \begin{pgfonlayer}{codebg}
    \fill[VSDarkBg, rounded corners=4pt] (-0.15,0.55) rectangle ({#2+0.45},{-\codepitch*#1-0.15});
    % the gutter with the line numbers, as in an editor
    \fill[white, opacity=0.05] (-0.15,0.12) rectangle (0.4,{-\codepitch*#1-0.06});
  \end{pgfonlayer}
  \foreach \d in {0,0.3,0.6} \fill[nLine] ({0.12+\d},0.33) circle (2.4pt);
  \coordinate (codeR) at ({#2+0.45},0);}
\newcommand{\cl}[3][VSText]{%
  \node[anchor=east, font=\ttfamily\tiny, text=VSText!45!VSDarkBg, inner sep=1pt] at (0.34,{-\codepitch*#2+0.14}) {#2};
  \node[anchor=west, font=\ttfamily\scriptsize, text=#1, inner sep=1pt] (L#2) at (0.5,{-\codepitch*#2+0.14}) {\strut #3};}
% the line that runs now, as in a debugger stopped at a breakpoint: the whole line lit from edge to edge,
% an accent on the left edge and an arrow in the gutter. The same mark on every slide that steps through code.
\newcommand{\clmark}[1]{%
  \begin{pgfonlayer}{codehl}
    \fill[AccentOrange!30!VSDarkBg] (-0.15,{-\codepitch*#1+0.14-0.18}) rectangle (codeR |- 0,{-\codepitch*#1+0.14+0.18});
    \fill[AccentOrange] (-0.15,{-\codepitch*#1+0.14-0.18}) rectangle (-0.08,{-\codepitch*#1+0.14+0.18});
  \end{pgfonlayer}
  \node[text=AccentOrange, font=\scriptsize, inner sep=0pt] at (0.04,{-\codepitch*#1+0.14}) {$\blacktriangleright$};}
% a line we talk about (not running): lit a little, from edge to edge, no frame
\newcommand{\clhl}[1]{%
  \begin{pgfonlayer}{codehl}
    \fill[white!14!VSDarkBg] (-0.15,{-\codepitch*#1+0.14-0.18}) rectangle (codeR |- 0,{-\codepitch*#1+0.14+0.18});
  \end{pgfonlayer}}
\newcommand{\clnote}[2]{%
  \clhl{#1}
  \draw[->, thick, AccentOrange, >=Stealth] ([xshift=0.6cm]codeR |- 0,{-\codepitch*#1+0.14}) -- ([xshift=0.08cm]codeR |- 0,{-\codepitch*#1+0.14});
  \node[anchor=west, font=\footnotesize, text=AccentOrange, align=left] at ([xshift=0.65cm]codeR |- 0,{-\codepitch*#1+0.14}) {#2};}
\newcommand{\kw}[1]{\textcolor{VSBlue}{\textbf{#1}}}
\newcommand{\fn}[1]{\textcolor{VSPurple}{#1}}
\newcommand{\cm}[1]{\textcolor{VSComment}{\textit{#1}}}
\newcommand{\st}[1]{\textcolor{VSString}{#1}}
\newcommand{\tb}{\phantom{xxxx}}


% the generator: a dark box with cogs, like the block cipher; \gennode{name}{position}{size: 1 large, 0 small}
\newcommand{\gennode}[3]{%
  \ifnum#3=1
    \node[hwmain, minimum width=2.4cm, minimum height=1.5cm] (#1) at (#2) {};
    \node[text=white, font=\LARGE] at ([yshift=0.22cm]#1.center) {\faCogs};
    \node[text=white, font=\small\bfseries] at ([yshift=-0.42cm]#1.center) {generator};
  \else
    \node[hwmain, minimum width=1.9cm, minimum height=1.05cm] (#1) at (#2) {};
    \node[text=white, font=\large] at ([yshift=0.16cm]#1.center) {\faCogs};
    \node[text=white, font=\scriptsize\bfseries] at ([yshift=-0.3cm]#1.center) {generator};
  \fi}
\tikzset{
  bigxor/.style={circle, draw=nInk, very thick, fill=white, inner sep=0pt, minimum size=0.8cm, font=\Large},
  smallbyte/.style={minimum width=0.5cm, minimum height=0.45cm},
}


% \bytes{x of the first cell}{y}{style}{values}{label above}
\newcommand{\bytes}[5]{%
  \foreach \v [count=\i from 0] in {#4} \node[#3, smallbyte] at ({#1+0.55*\i},#2) {\v};
  \node[font=\scriptsize\bfseries, text=nInk] at ({#1+0.55},{#2+0.45}) {#5};}

% a small dark box with one or more lines of code
\newcommand{\pycode}[1]{%
  \begin{beamercolorbox}[rounded=true,sep=1.6mm,wd=\linewidth]{vscodelisting}
    \ttfamily\scriptsize\color{VSText}#1
  \end{beamercolorbox}}


% an invoice on paper: \invoice{name}{position}{amount}
\newcommand{\invoice}[3]{%
  \node[draw=nLine, thick, fill=white, minimum width=2.1cm, minimum height=2.5cm, rounded corners=1pt] (#1) at (#2) {};
  \node[font=\tiny\bfseries, text=nInk, anchor=north] at ([yshift=-0.12cm]#1.north) {INVOICE};
  \foreach \d in {0.6,0.8,1.0} \draw[nLight, thick] ([shift={(0.25,-\d)}]#1.north west) -- ([shift={(-0.25,-\d)}]#1.north east);
  \node[font=\tiny, text=nInk, anchor=south west] at ([shift={(0.15,0.55)}]#1.south west) {to pay:};
  \node[font=\ttfamily\scriptsize\bfseries, text=nInk, anchor=south] at ([yshift=0.15cm]#1.south) {#3};}


\title{Cryptography}
\subtitle{Salsa20 and ChaCha20}
\author{Ksawery Możdżyński}
\date{}

\begin{document}

% ======================================================================
% Colour mapping for the 4 by 4 state (added in Lecture 2)
%   constants : neutral (nFill), the label says "const"
%   key       : S4 amber, light fill (as every key)
%   nonce     : white with a full amber frame (public, as the IV in Lecture 1)
%   counter   : white with a dashed amber frame (public, changes by itself)
% ======================================================================
\tikzset{
  wcell/.style={draw=nLine, thick, fill=white, text=nInk, font=\ttfamily\scriptsize, minimum width=1.55cm,
    minimum height=0.55cm, inner sep=1pt},
  st-c/.style={wcell, fill=nFill},
  st-k/.style={wcell, fill=sAmber!25, draw=sAmber},
  st-n/.style={wcell, draw=sAmber, line width=1pt},
  st-t/.style={wcell, draw=sAmber, line width=1pt, dashed},
  st-x/.style={wcell},
}
% a cell that changed or matters now: a light orange fill and a thin orange frame (as in Lecture 1)
\tikzset{cellmark/.style={draw=AccentOrange, thick, fill=AccentOrange, fill opacity=0.18, draw opacity=1, rounded corners=1.5pt}}
% the 4 by 4 state: \statefig[cell width]{value/kind, ... (16 entries)}{extra TikZ}; cells are named W0 ... W15
\newcommand{\statefig}[3][1.65]{%
  \begin{tikzpicture}[>=Stealth]
    \edef\statelist{#2}
    \foreach \v/\s [count=\i from 0] in \statelist {
      \pgfmathtruncatemacro{\r}{\i/4}\pgfmathtruncatemacro{\c}{mod(\i,4)}
      \node[st-\s, minimum width=#1cm] (W\i) at ({(#1+0.08)*\c},{-0.65*\r}) {\v};
    }
    #3
  \end{tikzpicture}}
\newcommand{\statelegend}{{\scriptsize
  \tikz[baseline=-0.6ex]\node[st-c, minimum width=0.5cm, minimum height=0.3cm] {};\ constant\quad
  \tikz[baseline=-0.6ex]\node[st-k, minimum width=0.5cm, minimum height=0.3cm] {};\ key\quad
  \tikz[baseline=-0.6ex]\node[st-t, minimum width=0.5cm, minimum height=0.3cm] {};\ counter\quad
  \tikz[baseline=-0.6ex]\node[st-n, minimum width=0.5cm, minimum height=0.3cm] {};\ nonce}}
\newcommand{\markW}[1]{\draw[cellmark] (W#1.south west) rectangle (W#1.north east);}
% a black box with three inputs (key, nonce, counter) and one block out; #1 = the name in the box
\newcommand{\threeinputs}[1]{%
  \begin{tikzpicture}[>=Stealth]
    \node[actorlight=sAmber, font=\scriptsize\bfseries, minimum width=2.3cm] (k) at (0,0.9) {\faKey\ key {\normalfont(32 bytes)}};
    \node[draw=sAmber, line width=1pt, fill=white, rounded corners=3pt, font=\scriptsize\bfseries, minimum width=2.3cm, inner sep=4pt] (n) at (0,0) {nonce};
    \node[draw=sAmber, line width=1pt, dashed, fill=white, rounded corners=3pt, font=\scriptsize\bfseries, minimum width=2.3cm, inner sep=4pt] (t) at (0,-0.9) {counter};
    \node[hwmain, minimum width=3.0cm, minimum height=2.3cm] (box) at (3.9,0) {};
    \node[text=white, font=\LARGE] at (3.9,0.35) {\faCogs};
    \node[text=white, font=\small\bfseries, align=center] at (3.9,-0.45) {#1};
    \draw[data] (k.east) -- (box.west |- k.east); \draw[data] (n.east) -- (box.west |- n.east); \draw[data] (t.east) -- (box.west |- t.east);
    \draw[data, line width=1.6pt] (box.east) -- (6.4,0);
    \node[ksb, minimum width=2.2cm, minimum height=0.6cm] (o) at (7.6,0) {64 bytes};
    \node[font=\scriptsize, text=nInk, anchor=north] at (o.south) {one block of keystream};
  \end{tikzpicture}}
% a horizontal bar chart row: \hbar{y}{label}{value}{scale}{text}
\newcommand{\hbarrow}[5]{%
  \fill[nLine] (0,{#1-0.2}) rectangle ({#3*#4},{#1+0.2});
  \node[font=\small\bfseries, text=nInk, anchor=east] at (-0.15,#1) {#2};
  \node[font=\small, text=nInk, anchor=west] at ({#3*#4+0.15},#1) {#5};}
% the ChaCha20 quarter round in a code window: #1 = the line that runs now (0 = none), #2 = width
\newcommand{\qrcode}[1]{%
  \begin{tikzpicture}
    \codebox{4}{5.2}
    \cl{1}{a += b;\ \ d \textasciicircum= a;\ \ d = rotl(d, 16)}
    \cl{2}{c += d;\ \ b \textasciicircum= c;\ \ b = rotl(b, 12)}
    \cl{3}{a += b;\ \ d \textasciicircum= a;\ \ d = rotl(d, 8)}
    \cl{4}{c += d;\ \ b \textasciicircum= c;\ \ b = rotl(b, 7)}
    \ifnum#1>0 \clmark{#1}\fi
  \end{tikzpicture}}
% four words a, b, c, d; #5 = the changed words, as a list of letters
\newcommand{\abcd}[5]{%
  \begin{tikzpicture}
    \foreach \l/\v [count=\i from 0] in {a/#1, b/#2, c/#3, d/#4} {
      \node[font=\small\bfseries, text=nInk] at ({2.0*\i},0.5) {\l};
      \node[wcell, minimum width=1.85cm, font=\ttfamily\small] (Q\l) at ({2.0*\i},0) {\v};
    }
    \foreach \l in {#5} {\draw[cellmark] (Q\l.south west) rectangle (Q\l.north east);}
  \end{tikzpicture}}

% ======================================================================
% Opening
% ======================================================================

\begin{frame}
  \titlepage
  \say{Dziś: szyfr strumieniowy, który działa w przeglądarce i w telefonie każdego z nas. Zobaczymy go od środka, linijka po linijce.}
\end{frame}

\begin{frame}{Agenda For Today}
  \vfill
  \begin{columns}[T,onlytextwidth]
    \column{0.52\textwidth}
      \large
      \begin{itemize}
        \item[\textcolor{DeepBlue}{\faStepForward}] \textbf{1. A keystream from a counter} \\
        \textcolor{gray}{\small Jump to any block}
        \vspace{0.3cm}

        \item[\textcolor{DeepBlue}{\faCalculator}] \textbf{2. Three simple operations} \\
        \textcolor{gray}{\small Add, rotate, XOR}
        \vspace{0.3cm}

        \item[\textcolor{DeepBlue}{\faTh}] \textbf{3. Salsa20} \\
        \textcolor{gray}{\small A 4 by 4 table of words}
        \vspace{0.3cm}

        \item[\textcolor{DeepBlue}{\faCogs}] \textbf{4. ChaCha20: one block} \\
        \textcolor{gray}{\small Twenty rounds, step by step}
      \end{itemize}
    \column{0.46\textwidth}
      \large
      \begin{itemize}
        \item[\textcolor{DeepBlue}{\faSearch}] \textbf{5. Why twenty rounds?} \\
        \textcolor{gray}{\small Measure it, then break it}
        \vspace{0.3cm}

        \item[\textcolor{DeepBlue}{\faRocket}] \textbf{6. Fast on any processor} \\
        \textcolor{gray}{\small No special hardware needed}
        \vspace{0.3cm}

        \item[\textcolor{DeepBlue}{\faIcon{shield-alt}}] \textbf{7. ChaCha20-Poly1305} \\
        \textcolor{gray}{\small Where you use it every day}
      \end{itemize}
  \end{columns}
  \vfill
  \say{Siedem części. Najpierw użyjemy szyfru jak czarnej skrzynki, potem zajrzymy do środka, a na końcu zobaczymy, gdzie działa naprawdę.}
\end{frame}

\begin{frame}{Today you will need}
  \vfill
  \begin{center}
    \begin{tikzpicture}
      \foreach \ic/\t/\s [count=\k from 0] in {\faPython/Python 3/{nothing to install, no extra libraries},
          \faEdit/An editor/{any one you like}, \faIcon{pencil-alt}/Paper and a pen/{for predictions before you run the code}} {
        \node[hw, minimum width=3.4cm, minimum height=2.4cm] (b\k) at ({4.2*\k},0) {};
        \node[text=nInk] at ({4.2*\k},0.35) {\fontsize{34}{36}\selectfont\ic};
        \node[font=\normalsize\bfseries, text=nInk] at ({4.2*\k},-0.65) {\t};
        \node[font=\scriptsize, text=gray, below=0.12cm of b\k, text width=3.6cm, align=center] {\s};
      }
    \end{tikzpicture}
  \end{center}
  \vfill
  \begin{center}\small The tasks:\ \ghlink{}\ \ Lecture2 / Task1 \ldots\ Task9\end{center}
  \vfill
  \say{Wszystkie zadania działają na czystym Pythonie, bez instalowania pakietów. Każde zadanie jest samodzielne: kto nie skończy jednego, może zacząć następne.}
\end{frame}

% ======================================================================
\section[A keystream from a counter]{A keystream from a counter}
% ======================================================================

\begin{frame}{Recall: writing bytes in hexadecimal}
  \begin{center}\small One byte = 8 bits = two hexadecimal digits (0 to 9, then A to F).\end{center}
  \vfill
  \begin{center}
    \begin{tikzpicture}[>=Stealth]
      \node[rowlabel, font=\small] at (-0.4,0) {binary};
      \foreach \b [count=\k from 0] in {0,1,0,0,1,0,0,0} {
        \pgfmathsetmacro{\x}{0.55*\k + (\k>3)*0.45}
        \node[byte, minimum width=0.48cm, font=\ttfamily\small] at (\x,0) {\b};
      }
      \draw[decorate, decoration={brace, amplitude=5pt}, nInk] (-0.25,0.35) -- (1.9,0.35)
        node[midway, above=6pt, font=\large\bfseries\ttfamily, text=nInk] {4};
      \draw[decorate, decoration={brace, amplitude=5pt}, nInk] (2.35,0.35) -- (4.5,0.35)
        node[midway, above=6pt, font=\large\bfseries\ttfamily, text=nInk] {8};
      \node[font=\small, text=nInk, anchor=west] at (5.0,0.5) {hex:\ \ written \texttt{\textbf{0x48}} or \texttt{\textbf{48}}};
      \node[font=\small, text=nInk, anchor=west] at (5.0,-0.1) {decimal:\ \ 72,\ \ the letter \texttt{H}};
    \end{tikzpicture}
  \end{center}
  \vfill
  \begin{center}\small A 32-bit number is 4 bytes, so 8 hex digits:\quad \texttt{0x61707865}\end{center}
  \vfill
  \say{Krótkie przypomnienie, bo dziś wszystko będzie w hex. Liczby 32-bitowe zobaczymy już za chwilę, to po prostu 4 bajty.}
\end{frame}

\samedot
\begin{frame}{Recall: XOR}
  \vfill
  \begin{columns}[c,onlytextwidth]
    \column{0.42\textwidth}
      \centering\Large
      \begin{tabular}{cc|c}
        \toprule
        $a$ & $b$ & $a$ XOR $b$ \\
        \midrule
        0 & 0 & 0 \\
        0 & 1 & \textbf{1} \\
        1 & 0 & \textbf{1} \\
        1 & 1 & 0 \\
        \bottomrule
      \end{tabular}
    \column{0.55\textwidth}
      \stepbox{Same bits give 0, different bits give 1.}{In Python:\ \texttt{a \textasciicircum{} b}}
      \stepbox{Twice with the same value changes nothing}{$(a \text{ XOR } k) \text{ XOR } k = a$}
  \end{columns}
  \vfill
  \keymsg{XOR with the same value twice gives back what we started with.}
  \say{Jeśli grupa pamięta XOR, wystarczy jedno zdanie. To ta sama operacja, której używamy do szyfrowania.}
\end{frame}

\samedot
\begin{frame}{Recall: a stream cipher}
  \vfill
  \begin{center}
    \begin{tikzpicture}[>=Stealth]
      \node[actorlight=sAmber, font=\small\bfseries, minimum height=0.8cm] (k) at (0,1.75) {\faKey\ key};
      \gennode{g}{3.2,1.75}{0}
      \draw[data] (k) -- (g);
      \draw[data] (g.east) -- (5.0,1.75);
      \byterow{5.4}{1.75}{ksb}{4A,91,0E,C7,5B}{0.62}
      \node[font=\scriptsize, text=nInk, anchor=west] at (8.4,1.75) {keystream \textbf{Z}\ \ (looks random)};
      \node[bigxor] (x) at (3.2,0) {$\oplus$};
      \draw[data] (g.south) -- (x.north);
      \node[actorlight=sViolet, font=\small\bfseries] (p) at (0,0) {plaintext \textbf{P}};
      \node[actorlight=sBrown, font=\small\bfseries] (c) at (6.6,0) {ciphertext \textbf{C}};
      \draw[data, line width=1.6pt] (p) -- (x); \draw[data, line width=1.6pt] (x) -- (c);
    \end{tikzpicture}
  \end{center}
  \vfill
  \begin{center}\large C = P XOR Z\qquad\qquad P = C XOR Z\end{center}
  \vfill
  \keymsg{Encryption and decryption are the same XOR with the same keystream.}
  \say{Generator rozwija krótki klucz w długi ciąg bajtów. Szyfrowanie to XOR bajt po bajcie. Odbiorca robi dokładnie to samo.}
\end{frame}

\samedot
\begin{frame}{Recall: never the same keystream twice}
  \begin{center}\small Two messages encrypted with the same keystream Z:\end{center}
  \vfill
  \begin{center}
    \begin{tikzpicture}[>=Stealth]
      \envelope{e1}{0,0.8}{C1}
      \envelope{e2}{0,-0.8}{C2}
      \node[actorlight=sAmber, font=\scriptsize\bfseries] at (1.45,1.25) {Z};
      \node[actorlight=sAmber, font=\scriptsize\bfseries] at (1.45,-0.35) {Z};
      \draw[data, line width=1.6pt] (1.5,0) -- (2.6,0);
      \node[font=\normalsize, text=nInk, anchor=west] at (2.7,0) {C1 XOR C2 = (P1 XOR Z) XOR (P2 XOR Z) = \textbf{P1 XOR P2}};
    \end{tikzpicture}
  \end{center}
  \vfill
  \begin{center}\small The keystream disappears. Knowing P1 gives P2.\\[3pt]
    So every message gets a public value, the \textbf{nonce}, that changes the keystream.\end{center}
  \vfill
  \keymsg{One key, many messages: each message needs its own nonce.}
  \say{Nonce to jawna wartość, inna dla każdej wiadomości. Klucz zostaje ten sam, a keystream jest inny. Dziś zobaczymy, jak nonce wchodzi do środka szyfru.}
\end{frame}

% a video player: #1 = 1: also the tape of encrypted bytes under it
\newcommand{\player}{%
  \draw[nInk, fill=nInk, rounded corners=4pt] (0,0) rectangle (11,3.0);
  \fill[nInk!80!black] (0.15,0.75) rectangle (10.85,2.85);
  \node[text=white!70, font=\fontsize{30}{30}\selectfont] at (5.5,1.8) {\faPlayCircle};
  \node[font=\scriptsize, text=white!80, anchor=north west] at (0.25,2.8) {encrypted film, 2 GB, 120 minutes};
  \fill[nLine] (0.4,0.38) rectangle (10.6,0.48);
  \fill[AccentOrange] (0.4,0.38) rectangle (8.05,0.48);
  \fill[white] (8.05,0.43) circle (0.12);
  \node[font=\tiny, text=white!80, anchor=north] at (0.4,0.33) {0 min};
  \node[font=\tiny, text=white!80, anchor=north] at (10.6,0.33) {120 min};}

\begin{frame}{A film on a server}
  \vfill
  \begin{center}
    \begin{tikzpicture}[>=Stealth]
      \player
      \draw[->, very thick, AccentOrange] (8.05,-0.6) -- (8.05,0.25);
      \node[font=\small\bfseries, text=AccentOrange, anchor=north] at (8.05,-0.6) {the viewer jumps here: 90 min};
    \end{tikzpicture}
  \end{center}
  \vfill
  \keymsg{To decrypt minute 90, we need the keystream for minute 90.}
  \say{Zwykła sytuacja: film na serwerze jest zaszyfrowany, widz przewija. Zapytać salę, co musi policzyć odtwarzacz.}
\end{frame}

% the keystream of the film as a tape: #1 = 1 mark the part to throw away, 2 cross it out
\newcommand{\filmtape}[1]{%
  \begin{tikzpicture}[>=Stealth]
    \fill[nIdle] (0,0) rectangle (8.25,0.5);
    \draw[nLine, thick] (0,0) rectangle (11,0.5);
    \fill[sAmber!45] (8.25,0) rectangle (8.9,0.5);
    \draw[sAmber, thick] (8.25,0) rectangle (8.9,0.5);
    \node[font=\scriptsize, text=nInk] at (4.1,0.25) {1.5 GB of keystream to throw away};
    \node[font=\scriptsize\bfseries, text=sAmberText, anchor=south] at (8.57,0.55) {needed};
    \node[font=\tiny, text=gray, anchor=north] at (0,-0.05) {0 min};
    \node[font=\tiny, text=gray, anchor=north] at (8.25,-0.05) {90 min};
    \node[font=\tiny, text=gray, anchor=north] at (11,-0.05) {120 min};
    \ifnum#1=2
      \draw[flow] (8.57,1.65) -- (8.57,0.95);
      \node[font=\scriptsize\bfseries, text=DeepBlue, anchor=south] at (8.57,1.65) {straight there};
    \fi
  \end{tikzpicture}}

\samedot
\begin{frame}{A generator that only goes forward}
  \begin{center}\small A generator like \textbf{RC4} gives the keystream one byte after another.\\
    To reach minute 90 it must first produce and throw away 1.5 GB of keystream.\end{center}
  \vfill
  \begin{center}\filmtape{1}\end{center}
  \vfill
  \begin{center}
    % CHECKED (draft, author's laptop): openssl speed, RC4 0.57 GB/s; 1.5 / 0.57 = 2.6 s
    \small
    \begin{tabular}{ll}
      measured RC4 speed on a laptop: & 0.57 GB/s \\
      time before the film can play: & about \textbf{2.6 s}, for every jump \\
    \end{tabular}
  \end{center}
  \vfill
  \say{RC4 to szyfr z poprzedniego wykładu: każdy jego bajt zależy od wszystkich wcześniejszych. 1,5 GB to trzy czwarte filmu. Prędkość RC4 zmierzyłem openssl speed na tym laptopie; przy każdym przewinięciu trzeba liczyć od początku.}
\end{frame}

\samedot
\begin{frame}{What do we really need?}
  \vfill
  \begin{center}\Large Any piece of the keystream,\\ without computing the pieces before it.\end{center}
  \vfill
  \begin{center}\filmtape{2}\end{center}
  \vfill
  \keymsg{We want to jump straight to any part of the keystream.}
  \say{Poczekać na pomysły z sali. Ktoś zwykle zaproponuje numerowanie kawałków.}
\end{frame}

\begin{frame}{The idea: number the blocks}
  \begin{center}\small Cut the keystream into blocks of \textbf{64 bytes} and give each block a number.\end{center}
  \vfill
  \begin{center}
    \begin{tikzpicture}[>=Stealth]
      \foreach \n/\x in {0/0, 1/1.9, 2/3.8} {
        \node[ksb, minimum width=1.7cm, minimum height=0.7cm] (b\n) at (\x,0) {64 bytes};
        \node[font=\scriptsize\bfseries, text=nInk, anchor=south] at (\x,0.4) {block \n};
      }
      \node[font=\large, text=nInk] at (5.4,0) {\ldots};
      \node[ksb, minimum width=1.7cm, minimum height=0.7cm] (bn) at (7.2,0) {64 bytes};
      \node[font=\scriptsize\bfseries, text=nInk, anchor=south] at (7.2,0.4) {block 23\,437\,500};
      \node[font=\large, text=nInk] at (8.7,0) {\ldots};
    \end{tikzpicture}
  \end{center}
  \vfill
  \begin{center}\small Block $n$ is computed from three things only:\quad
    \tikz[baseline=-0.6ex]\node[actorlight=sAmber, font=\scriptsize\bfseries, inner sep=2pt] {key};,\
    \tikz[baseline=-0.6ex]\node[draw=sAmber, line width=1pt, rounded corners=3pt, font=\scriptsize\bfseries, inner sep=2pt] {nonce};,\
    \tikz[baseline=-0.6ex]\node[draw=sAmber, line width=1pt, dashed, rounded corners=3pt, font=\scriptsize\bfseries, inner sep=2pt] {$n$};\end{center}
  \vfill
  \keymsg{Block $n$ does not depend on blocks 0 to $n - 1$.}
  % CHECKED: 1.5e9 / 64 = 23 437 500
  \say{1,5 GB to 23 437 500 bloków po 64 bajty. Ten numer liczymy jednym dzieleniem, a nie liczeniem bloków po kolei.}
\end{frame}

\samedot
\begin{frame}{The counter}
  \bigdef{\textbf{Counter}: the number of a 64-byte block of the keystream. It goes 0, 1, 2, \ldots\
    for the blocks of one message. Like the nonce, it is public.}
  \vfill
  \begin{center}
    \small
    \begin{tabular}{l|c|c|c|c}
      message: & bytes 0 to 63 & bytes 64 to 127 & bytes 128 to 191 & \ldots \\
      \hline
      counter: & 0 & 1 & 2 & \ldots \\
    \end{tabular}
  \end{center}
  \vfill
  \begin{center}\small Byte number \texttt{b} of the message is in block\ \texttt{\textbf{b // 64}}\end{center}
  \vfill
  \say{Licznik nie jest tajny: odbiorca wylicza go sam z pozycji bajtu. Tajny jest tylko klucz.}
\end{frame}

\samedot
\begin{frame}{Salsa20 and ChaCha20 from the outside}
  \vfill
  \begin{center}\threeinputs{Salsa20 /\\ChaCha20}\end{center}
  \vfill
  \begin{center}\small Two ciphers built on this idea, by Daniel J. Bernstein:\\
    \textbf{Salsa20} (2005) and \textbf{ChaCha20} (2008), an improved Salsa20.\end{center}
  \vfill
  \keymsg{Same key and nonce, a new counter: a new block.}
  \say{Na razie traktujemy szyfr jak czarną skrzynkę. ChaCha20 jest następcą Salsa20, o różnicach powiemy w sekcji 3 i 4. Liczba 20 w nazwie to liczba rund, o tym w sekcji 5.}
\end{frame}

\begin{frame}{Task 1: use ChaCha20}
  \taskhead{Task 1}{You received an encrypted message and you know the key and the nonce.}
  \vfill
  \begin{center}
    \small
    \renewcommand{\arraystretch}{1.5}
    \begin{tabular}{l l}
      key & \texttt{c0ffee00 c0ffee00 \ldots} (32 bytes) \\
      nonce & \texttt{0000004a 00000000 00000007} \\
      ciphertext & \texttt{36323fcd0cbf7202049805976e0d4afda0a2bc363f02eea05b\ldots} \\
    \end{tabular}
  \end{center}
  \vfill
  \begin{center}\large The message was encrypted starting with \textbf{counter 1}.\\[3pt] \textbf{What does it say?}\end{center}
  \vfill
  \say{Około 4 minut. Pierwsze zadanie tylko z użyciem gotowej funkcji, bez zaglądania do środka.}
\end{frame}

\samedot
\begin{frame}[fragile]{Task 1: what to do}
  \begin{columns}[T,onlytextwidth]
    \column{0.6\textwidth}
\begin{lstlisting}[style=py, emph={chacha20,print}, basicstyle=\color{VSText}\ttfamily\fontsize{6.4}{7.4}\selectfont]
from chacha import chacha20

key = bytes.fromhex("c0ffee00" * 8)              # 32 bytes
nonce = bytes.fromhex("0000004a0000000000000007")   # 12 bytes
ciphertext = bytes.fromhex(
    "36323fcd0cbf7202049805976e0d4afda0a2bc36"
    "3f02eea05b49125427ad01f9709e3275ef98d8"
)

# TODO 1: decrypt the ciphertext with
#         chacha20(key, nonce, counter, data).
#         The message was encrypted starting with counter 1.

# TODO 2: print the plaintext as text: plaintext.decode()
\end{lstlisting}
    \column{0.39\textwidth}
      \taskhead{Task 1}{Open \texttt{use\_chacha.py} from \ghlink{Task1}.}
      \footnotesize
      \begin{enumerate}\setlength{\itemsep}{4pt}
        \item Call \texttt{chacha20} with counter 1.
        \item Print the result as text.
      \end{enumerate}
  \end{columns}
  \say{chacha.py jest w folderze zadania, nie trzeba go czytać. Szyfrowanie i odszyfrowanie to ta sama funkcja.}
\end{frame}

\samedot
\begin{frame}[fragile]{Solution: use ChaCha20}
  \begin{columns}[T,onlytextwidth]
    \column{0.55\textwidth}
% CHECKED: Task1/solution/use_chacha.py prints "The meeting moves to room 204 at 15:30."
\begin{lstlisting}[style=py, emph={chacha20,print}]
plaintext = chacha20(key, nonce, 1, ciphertext)
print(plaintext.decode())
\end{lstlisting}
    \column{0.42\textwidth}
      \stepbox{The output}{\texttt{The meeting moves to room 204 at 15:30.}}
  \end{columns}
  \vfill
  \keymsg{To decrypt, call the same function again.}
  \say{Kod sprawdzony, daje dokładnie ten tekst. Pytanie do sali: co zmienić, żeby odszyfrować tylko koniec wiadomości?}
\end{frame}

\begin{frame}{Task 2: jump to the middle}
  \taskhead{Task 2}{\texttt{film.bin}: an encrypted film of 1000 blocks, 64 bytes each.
    It was encrypted starting with counter 0: block $n$ used counter $n$.}
  \vfill
  \begin{center}
    \begin{tikzpicture}
      \foreach \k in {0,...,39} {
        \ifnum\k=28 \node[ctb, minimum width=0.22cm, minimum height=0.6cm, inner sep=0pt] at ({0.26*\k},0) {};
          \draw[annot] ({0.26*\k-0.15},-0.37) rectangle ({0.26*\k+0.15},0.37);
        \else \node[ctb, minimum width=0.22cm, minimum height=0.6cm, inner sep=0pt, fill=sBrown!8] at ({0.26*\k},0) {}; \fi
      }
      \node[font=\scriptsize\bfseries, text=AccentOrange, anchor=north] at ({0.26*28},-0.4) {block 700};
      \node[font=\tiny, text=gray, anchor=north] at (0,-0.4) {block 0};
      \node[font=\tiny, text=gray, anchor=north] at ({0.26*39},-0.4) {block 999};
    \end{tikzpicture}
  \end{center}
  \vfill
  \small The program already decrypts the whole film and prints block 700.\\[3pt]
  \textbf{Decrypt only block 700}, without blocks 0 to 699, and compare the times.
  \vfill
  \say{Około 5 minut. To samo co w zadaniu 1, tylko trzeba pomyśleć o liczniku. (Rysunek pokazuje film w skrócie, nie wszystkie 1000 bloków.)}
\end{frame}

\samedot
\begin{frame}[fragile]{Task 2: what to do}
  \begin{columns}[T,onlytextwidth]
    \column{0.6\textwidth}
\begin{lstlisting}[style=py, emph={chacha20,print,open}, basicstyle=\color{VSText}\ttfamily\fontsize{6.4}{7.4}\selectfont]
film = open("film.bin", "rb").read()       # 1000 blocks of 64 bytes
# Encrypted starting with counter 0: block n used counter n.

start = time.perf_counter()
everything = chacha20(key, nonce, 0, film)
print("whole film:", everything[700 * 64:701 * 64].decode().strip())
print(f"  took {1000 * (time.perf_counter() - start):.1f} ms")

start = time.perf_counter()
# TODO 1: cut block 700 out of the film
#         (bytes 700*64 to 701*64 of the film).

# TODO 2: decrypt ONLY that block.
#         Which counter do you need?
one_block = b""
\end{lstlisting}
      {\scriptsize\textcolor{gray}{Shortened on the slide: the imports, the key, the nonce and the last two prints are in the file.}}
    \column{0.37\textwidth}
      \taskhead{Task 2}{Open \texttt{jump.py} from \ghlink{Task2}.}
      \footnotesize
      \begin{enumerate}\setlength{\itemsep}{4pt}
        \item Cut block 700 out of the film.
        \item Decrypt it with the right counter.
      \end{enumerate}
      \stepbox{In Python}{a slice of bytes:\ \texttt{film[start:end]}}
  \end{columns}
  \say{Pomiar czasu jest już w pliku. Fragment bajtów w Pythonie: film[od:do].}
\end{frame}

\samedot
\begin{frame}[fragile]{Solution: jump to the middle}
% CHECKED: Task2/solution/jump.py, author's laptop: 135 to 138 ms and 0.1 to 0.2 ms in three runs
\begin{lstlisting}[style=py, emph={chacha20}]
piece = film[700 * 64:701 * 64]
one_block = chacha20(key, nonce, 700, piece)
\end{lstlisting}
  \vfill
  \begin{center}
    \begin{tikzpicture}
      \hbarrow{0.6}{whole film}{136}{0.05}{136 ms}
      \hbarrow{0}{one block}{0.1}{0.05}{0.1 ms}
    \end{tikzpicture}\\[4pt]
    \small both print:\ \texttt{Frame 0700: the hero finally opens the red door.}
  \end{center}
  \vfill
  \keymsg{The counter lets us decrypt any block on its own.}
  \say{Kod sprawdzony; czasy z mojego laptopa, trzy uruchomienia dały 135 do 138 ms i 0,1 do 0,2 ms. U studentów liczby będą inne, ale różnica podobna.}
\end{frame}

\begin{frame}{Nonce and counter: two different jobs}
  \vfill
  \begin{center}
    \small
    \renewcommand{\arraystretch}{1.6}
    \begin{tabular}{l l l}
      \toprule
      & changes between \ldots & who chooses it \\
      \midrule
      \tikz[baseline=-0.6ex]\node[draw=sAmber, line width=1pt, rounded corners=3pt, font=\scriptsize\bfseries, inner sep=3pt] {nonce}; & messages sent with one key & the sender, once per message \\
      \tikz[baseline=-0.6ex]\node[draw=sAmber, line width=1pt, dashed, rounded corners=3pt, font=\scriptsize\bfseries, inner sep=3pt] {counter}; & blocks of one message & counts 0, 1, 2, \ldots\ by itself \\
      \bottomrule
    \end{tabular}
  \end{center}
  \vfill
  \begin{center}\small Both are public. The \textbf{key} is the only secret.\end{center}
  \vfill
  \keymsg{A (nonce, counter) pair must never repeat with the same key.}
  \say{Powtórka pary nonce i licznik to dokładnie powtórka keystreamu, czyli błąd, który usuwa klucz.}
\end{frame}

\samedot
\begin{frame}{How long can one message be?}
  \begin{center}\small In ChaCha20 as used on the internet (\textbf{RFC 8439}):\end{center}
  \vfill
  \begin{center}
    \begin{tikzpicture}[>=Stealth]
      \foreach \k in {0,...,31} \node[draw=sAmber, dashed, thick, fill=white, minimum width=0.26cm, minimum height=0.4cm, inner sep=0pt] at ({0.28*\k},0) {};
      \node[font=\scriptsize\bfseries, text=nInk, anchor=south] at (4.34,0.3) {counter: 32 bits $\rightarrow$ $2^{32}$ blocks};
      \draw[data] (9.2,0) -- (10.2,0);
      \node[font=\small, text=nInk, anchor=west, align=left] at (10.3,0) {$2^{32} \times 64$ bytes\\ $=$ \textbf{256 GiB}};
    \end{tikzpicture}
  \end{center}
  \vfill
  \begin{center}\small for one key and nonce\end{center}
  \vfill
  \keymsg{After 256 GiB, a new nonce is needed.}
  % CHECKED: 2**32 * 64 / 2**30 = 256
  \say{RFC to dokument z oficjalną specyfikacją standardu internetowego. Większych wiadomości w praktyce się nie wysyła; protokoły internetowe i tak dzielą dane na małe kawałki.}
\end{frame}

% ======================================================================
\section[Three simple operations]{Three simple operations}
% ======================================================================

\begin{frame}{Inside: only 32-bit words}
  \begin{center}\small ChaCha20 does not work on single bytes. It works on \textbf{words}.\end{center}
  \vfill
  \begin{columns}[c,onlytextwidth]
    \column{0.45\textwidth}
      \centering
      \begin{tikzpicture}
        \foreach \b [count=\k from 0] in {61,70,78,65} \node[byte, minimum width=0.8cm, font=\ttfamily\small] at ({0.85*\k},0) {\b};
        \draw[decorate, decoration={brace, mirror, amplitude=5pt}, nInk] (-0.4,-0.35) -- (2.95,-0.35)
          node[midway, below=6pt, font=\scriptsize, text=nInk] {one word = 32 bits = 4 bytes};
        \node[font=\ttfamily\small\bfseries, text=nInk] at (1.27,0.65) {0x61707865};
      \end{tikzpicture}
    \column{0.5\textwidth}
      \centering
      \begin{tikzpicture}
        \foreach \i in {0,...,15} {\pgfmathtruncatemacro{\r}{\i/4}\pgfmathtruncatemacro{\c}{mod(\i,4)}
          \node[cell, minimum width=0.9cm, minimum height=0.45cm] at ({0.95*\c},{-0.5*\r}) {};}
        \node[font=\scriptsize, text=nInk, anchor=north] at (1.42,-1.85) {the state = 16 words = 64 bytes};
      \end{tikzpicture}
  \end{columns}
  \vfill
  \begin{center}\small Everything inside happens on these 16 words.\end{center}
  \vfill
  \say{Stan to cała pamięć szyfru w trakcie liczenia jednego bloku. 64 bajty stanu dają 64 bajty keystreamu.}
\end{frame}

\begin{frame}{Recall: mod}
  \vfill
  \begin{columns}[c,onlytextwidth]
    \column{0.45\textwidth}
      \centering
      \begin{tikzpicture}[>=Stealth]
        \draw[nLine, very thick, fill=nFill] (0,0) circle (1.6);
        \foreach \h in {1,...,12} \node[font=\small, text=nInk] at ({90-30*\h}:1.28) {\h};
        \draw[flow] (-60:1.95) arc (-60:-210:1.95);
        \node[font=\scriptsize\bfseries, text=DeepBlue] at (0,0) {+5 hours};
      \end{tikzpicture}
    \column{0.52\textwidth}
      \stepbox{$a$ mod $n$}{the remainder after dividing $a$ by $n$.\\ \texttt{7 mod 4 = 3},\quad \texttt{8 mod 4 = 0}}
      \stepbox{A clock}{10 o'clock + 5 hours = 3 o'clock\\ \texttt{(10 + 5) mod 12 = 3}}
      \stepbox{In Python}{\texttt{7 \% 4}}
  \end{columns}
  \vfill
  \say{Mod przyda się za chwilę: dodawanie w ChaCha20 zawija się jak zegar, tylko tarcza ma 2 do potęgi 32 pozycji.}
\end{frame}

\samedot
\begin{frame}{Operation 1: addition mod $2^{32}$}
  \begin{center}\small A word has 32 bits. A carry out of bit 31 is \textbf{lost}.\end{center}
  \vfill
  \begin{center}
    \ttfamily\small
    \begin{tabular}{r r l}
      & 0xFFFFFFFF & \quad 1111 1111 \ldots\ 1111 \\
      + & 0x00000001 & \quad 0000 0000 \ldots\ 0001 \\
      \hline
      = & 0x00000000 & \tikz[baseline=-0.6ex]{\node[inner sep=1pt, text=AccentOrange, font=\ttfamily\small] (o) {1}; \draw[AccentOrange, thick] (o.south west) -- (o.north east);}\ 0000 0000 \ldots\ 0000 \\
    \end{tabular}\\[4pt]
    {\normalfont\scriptsize\textcolor{AccentOrange}{the 33rd bit is dropped}}
  \end{center}
  \vfill
  \begin{center}\small It is addition mod $2^{32}$: the clock has 4\,294\,967\,296 positions.\end{center}
  \vfill
  \say{To samo robi procesor przy zwykłym dodawaniu liczb 32-bitowych, więc ta operacja jest bardzo szybka.}
\end{frame}

\samedot
\begin{frame}{Addition in Python}
  \begin{center}\small Python integers have no limit, so we cut the result to 32 bits ourselves:\end{center}
  \pycode{\large (a + b) \& 0xFFFFFFFF}
  \vfill
  \begin{center}\small \texttt{0xFFFFFFFF} is 32 ones. \texttt{\&} keeps only the lowest 32 bits.\end{center}
  \vfill
  \begin{center}
    \ttfamily\small
    \begin{tabular}{l l}
      0xFFFFFFFF + 1 & = 0x\textcolor{AccentOrange}{\textbf{1}}00000000 \\
      (0xFFFFFFFF + 1) \& 0xFFFFFFFF & = 0x00000000 \\
    \end{tabular}
  \end{center}
  \vfill
  \say{\& to bitowe AND: bit zostaje tylko tam, gdzie maska ma 1. W C lub w procesorze nie trzeba tego robić, ucięcie dzieje się samo.}
\end{frame}

% a row of 32 bits in groups of 4: \bitrow{y}{label}{bits as a string of 8 groups separated by spaces}
\begin{frame}{Operation 2: XOR on words}
  \begin{center}\small The same XOR as for bytes, on 32 bits at once:\end{center}
  \vfill
  \begin{center}
    % CHECKED: 0x12131415 ^ 0x01234567 = 0x13305172
    \ttfamily\small
    \renewcommand{\arraystretch}{1.4}
    \begin{tabular}{r r l}
      & 0x12131415 & \quad 0001 0010 0001 0011 0001 0100 0001 0101 \\
      XOR & 0x01234567 & \quad 0000 0001 0010 0011 0100 0101 0110 0111 \\
      \hline
      = & 0x13305172 & \quad 0001 0011 0011 0000 0101 0001 0111 0010 \\
    \end{tabular}
  \end{center}
  \vfill
  \say{Te liczby nie są przypadkowe: zobaczymy je znowu w sekcji 4.}
\end{frame}

% 32 bits as cells: \bitcells[name]{y}{bit string of 32}{label}; cells are named <name>-0 ... <name>-31
\newcommand{\bitcells}[4][B]{%
  \node[rowlabel, font=\ttfamily\small] at (-0.3,#2) {#4};
  \foreach \k in {0,...,31} {
    \pgfmathsetmacro{\bx}{0.3*\k + 0.12*floor(\k/4)}
    \node[cell, minimum width=0.27cm, minimum height=0.36cm, inner sep=0pt, font=\ttfamily\tiny] (#1-\k) at (\bx,#2) {\StrChar{#3}{\the\numexpr\k+1\relax}};
  }}

\begin{frame}{Operation 3: rotation}
  \defbox{\textbf{Rotate left} by $n$: the bits that fall out on the left come back on the right.}
  \vfill
  \begin{center}
    \begin{tikzpicture}[>=Stealth]
      \bitcells[X]{0}{10000000000000000000000000000001}{x}
      \node[font=\ttfamily\small, text=nInk, anchor=west] at (10.6,0) {0x80000001};
      \bitcells[R]{-1.6}{00000000000000000000000000000011}{rotl(x, 1)}
      \node[font=\ttfamily\small, text=nInk, anchor=west] at (10.6,-1.6) {0x00000003};
      \draw[cellmark] (X-0.south west) rectangle (X-0.north east);
      \draw[cellmark] (R-31.south west) rectangle (R-31.north east);
      \draw[->, thick, AccentOrange] (X-0.south) -- (0,-0.8) -- ({0.3*31+0.12*7},-0.8) -- (R-31.north);
      \node[font=\scriptsize\bfseries, text=AccentOrange, fill=white, inner sep=1pt] at (5.0,-0.8) {the top bit comes back at the bottom};
    \end{tikzpicture}
  \end{center}
  \vfill
  \begin{center}\small No bit is lost, the word just turns around.\end{center}
  \vfill
  \keymsg{Rotation moves bits to new positions without losing any.}
  \say{Rotacja przenosi bity na inne pozycje, dzięki czemu bity z góry słowa wpływają na dół i odwrotnie. Pozostałe bity przesuwają się o jedno miejsce w lewo.}
\end{frame}

\samedot
\begin{frame}{Rotation is not a shift}
  \vfill
  \begin{center}
    \begin{tikzpicture}[>=Stealth]
      \bitcells[X]{0}{10000000000000000000000000000001}{x}
      \bitcells[S]{-0.9}{00000000000000000000000000000010}{x << 1}
      \node[font=\scriptsize, text=nInk, anchor=west] at (10.5,-0.9) {\nok\ top bit lost};
      \bitcells[T]{-1.8}{00000000000000000000000000000011}{rotl(x, 1)}
      \node[font=\scriptsize, text=nInk, anchor=west] at (10.5,-1.8) {\ok\ top bit back};
      \draw[cellmark] (T-31.south west) rectangle (T-31.north east);
    \end{tikzpicture}
  \end{center}
  \vfill
  \begin{center}\small In Python, a rotation by $n$ is built from two shifts:\end{center}
  \pycode{\large ((x << n) | (x >> (32 - n))) \& 0xFFFFFFFF}
  \vfill
  \say{| to bitowe OR. Lewe przesunięcie robi miejsce, prawe przynosi bity z góry na dół. Wiele procesorów ma rotację jako jedną instrukcję.}
\end{frame}

\begin{frame}{ARX}
  \vfill
  \begin{center}
    \begin{tikzpicture}
      \foreach \ic/\t/\k in {{\textbf{+}}/Add/0, {\faRedo}/Rotate/1, {$\oplus$}/XOR/2} {
        \node[hw, minimum width=3.2cm, minimum height=2.4cm] at ({4.0*\k},0) {};
        \node[text=nInk, font=\fontsize{34}{36}\selectfont] at ({4.0*\k},0.3) {\ic};
        \node[font=\large\bfseries, text=nInk] at ({4.0*\k},-0.75) {\t};
      }
      \node[font=\huge\bfseries, text=DeepBlue] at (0,-1.6) {A};
      \node[font=\huge\bfseries, text=DeepBlue] at (4,-1.6) {R};
      \node[font=\huge\bfseries, text=DeepBlue] at (8,-1.6) {X};
    \end{tikzpicture}
  \end{center}
  \vfill
  \begin{center}\small Three operations that every processor does in one step.\\ No tables in memory, no special instructions.\end{center}
  \keymsg{Salsa20 and ChaCha20 are only these three operations, repeated many times.}
  \say{Szyfry z tej rodziny nazywa się ARX od pierwszych liter. Brak tablic okaże się ważny w sekcji 6.}
\end{frame}

\begin{frame}{Task 3: write the operations}
  \taskhead{Task 3}{Write two functions on 32-bit words. The program checks six examples and prints OK or FAIL for each.}
  \vfill
  \begin{columns}[T,onlytextwidth]
    \column{0.48\textwidth}
      \stepbox{\texttt{add32(a, b)}}{addition mod $2^{32}$\\[2pt] \texttt{add32(0xFFFFFFFF, 0x00000001)} $\rightarrow$ \texttt{0x00000000}}
    \column{0.48\textwidth}
      \stepbox{\texttt{rotl32(x, n)}}{rotation left by $n$ bits\\[2pt] \texttt{rotl32(0x80000001, 1)} $\rightarrow$ \texttt{0x00000003}}
  \end{columns}
  \vfill
  \say{Około 6 minut. Oba wzory były na poprzednich slajdach, ale niech spróbują bez patrzenia.}
\end{frame}

\samedot
\begin{frame}[fragile]{Task 3: what to do}
  \begin{columns}[T,onlytextwidth]
    \column{0.6\textwidth}
\begin{lstlisting}[style=py]
MASK = 0xFFFFFFFF   # 32 ones: keeps only the lowest 32 bits


def add32(a, b):
    # TODO 1: add a and b and keep only 32 bits
    pass


def rotl32(x, n):
    # TODO 2: rotate the 32-bit word x left by n bits
    pass
\end{lstlisting}
      {\scriptsize\textcolor{gray}{Shortened on the slide: the six tests and longer comments are in the file.}}
    \column{0.37\textwidth}
      \taskhead{Task 3}{Open \texttt{arx.py} from \ghlink{Task3}.}
      \stepbox{One line of the output}{\scriptsize\texttt{OK \ rotl32: got 0x00000003,\\ \ \ \ \ expected 0x00000003}}
      \footnotesize At the start all six tests give \texttt{FAIL}.
  \end{columns}
  \say{Testy są gotowe na dole pliku. Na początku wszystkie sześć daje FAIL.}
\end{frame}

\samedot
\begin{frame}[fragile]{Solution: write the operations}
  \begin{columns}[T,onlytextwidth]
    \column{0.58\textwidth}
% CHECKED: Task3/solution/arx.py prints OK six times
\begin{lstlisting}[style=py]
def add32(a, b):
    return (a + b) & MASK

def rotl32(x, n):
    return ((x << n) | (x >> (32 - n))) & MASK
\end{lstlisting}
    \column{0.39\textwidth}
      \stepbox{The output}{\texttt{OK} six times}
  \end{columns}
  \vfill
  \say{Kod sprawdzony. Te dwie funkcje są gotowe we wszystkich kolejnych zadaniach, więc nikt nie musi ich przepisywać.}
\end{frame}

\begin{frame}{Why mix all three?}
  \vfill
  \begin{columns}[T,onlytextwidth]
    \column{0.48\textwidth}
      \stepbox{XOR and rotation alone}{each output bit = XOR of some input bits\\[3pt]
        \texttt{\scriptsize y5 = x2 \textasciicircum{} x9 \textasciicircum{} x30}\\[3pt]
        the attacker can solve it like a set of equations}
    \column{0.48\textwidth}
      \stepbox{Addition}{the carry depends on two bits at once\\[3pt]
        \texttt{\scriptsize carry = x \& y}\quad (1 only if both are 1)\\[3pt]
        this is \textbf{not} a XOR of bits any more}
  \end{columns}
  \vfill
  \keymsg{Addition breaks the equations. XOR and rotation spread the bits.}
  \say{Zależność, którą da się zapisać jako XOR bitów, nazywamy liniową. Wrócimy do tego w sekcji 5. (Równanie y5 na slajdzie to tylko przykład postaci, nie konkretny bit ChaCha20.)}
\end{frame}

% ======================================================================
\section[Salsa20]{Salsa20}
% ======================================================================

\newcommand{\positions}{0/x,1/x,2/x,3/x,4/x,5/x,6/x,7/x,8/x,9/x,10/x,11/x,12/x,13/x,14/x,15/x}

\begin{frame}{The state: a 4 by 4 table}
  \begin{center}\small The 16 words of the state are drawn as a 4 by 4 table:\end{center}
  \vfill
  \begin{center}\statefig[1.9]{\positions}{}\end{center}
  \vfill
  \begin{center}\small The numbers are the \textbf{positions} of the words, 0 to 15.\end{center}
  \vfill
  \say{Ten sam układ ma Salsa20 i ChaCha20. Różni się tylko to, co wpisujemy w które pole.}
\end{frame}

\samedot
\begin{frame}{Salsa20: what goes where}
  \begin{columns}[c,onlytextwidth]
    \column{0.55\textwidth}
      \centering
      \statefig[1.6]{const/c,key/k,key/k,key/k,key/k,const/c,nonce/n,nonce/n,count/t,count/t,const/c,key/k,key/k,key/k,key/k,const/c}{}\\[6pt]
      \statelegend
    \column{0.42\textwidth}
      \small
      \renewcommand{\arraystretch}{1.4}
      \begin{tabular}{l r}
        key & 8 words = 256 bits \\
        nonce & 2 words = 64 bits \\
        counter & 2 words = 64 bits \\
        const & 4 words on the diagonal \\
      \end{tabular}
  \end{columns}
  \vfill
  \say{Klucz zajmuje połowę stanu. Napastnik zna stałe, nonce i licznik, nie zna tylko 8 słów klucza.}
\end{frame}

\samedot
\begin{frame}{The constants}
  \begin{center}\small The four constant words are ASCII text:\end{center}
  \vfill
  \begin{center}
    % CHECKED: python3, struct.pack("<4I", ...) of the four words = b"expand 32-byte k"
    \begin{tikzpicture}
      \foreach \w/\t [count=\k from 0] in {61707865/expa, 3320646e/{nd 3}, 79622d32/{2-by}, 6b206574/{te k}} {
        \node[st-c, minimum width=2.4cm, font=\ttfamily\small] at ({2.7*\k},0) {0x\w};
        \node[font=\ttfamily\normalsize, text=nInk] at ({2.7*\k},-0.65) {"\t"};
      }
    \end{tikzpicture}
  \end{center}
  \vfill
  \begin{center}\Large\texttt{"expand 32-byte k"}\end{center}
  \vfill
  \keymsg{A fixed part of the state that nobody can choose.}
  \say{Bajty w słowie czytamy od końca (wyjaśnimy to w sekcji 4), stąd 0x61707865 to expa. Stałe gwarantują, że nawet przy kluczu z samych zer stan nie jest pusty, i że napastnik nie wybiera całego stanu.}
\end{frame}

% the Salsa20 quarter round in a code window: #1 = the line that runs now (0 = none)
\newcommand{\salsacode}[1]{%
  \begin{tikzpicture}
    \codebox{4}{4.2}
    \cl{1}{b \textasciicircum= rotl(a + d, 7)}
    \cl{2}{c \textasciicircum= rotl(b + a, 9)}
    \cl{3}{d \textasciicircum= rotl(c + b, 13)}
    \cl{4}{a \textasciicircum= rotl(d + c, 18)}
    \ifnum#1>0 \clmark{#1}\fi
  \end{tikzpicture}}

\begin{frame}{The Salsa20 quarter round}
  \begin{columns}[c,onlytextwidth]
    \column{0.45\textwidth}
      \centering
      {\small quarter round of four words a, b, c, d}\\[4pt]
      \salsacode{0}
    \column{0.52\textwidth}
      \stepbox{Every line}{add two words, rotate, XOR into the third.}
      \stepbox{\texttt{+}}{is addition mod $2^{32}$.}
      \defbox{A \textbf{quarter round} mixes 4 of the 16 words.}
  \end{columns}
  \vfill
  \keymsg{The quarter round mixes four words with ARX.}
  \say{Quarter round to podstawowy klocek: działa na czterech słowach z szesnastu, stąd ćwierć. Runda Salsa20 to cztery quarter round.}
\end{frame}

\begin{frame}{Task 4: one quarter round of Salsa20}
  \taskhead{Task 4}{Input:\ \texttt{a = 0x00000001,\ b = 0,\ c = 0,\ d = 0}.
    Write the four lines and print the words after each line.}
  \vfill
  \begin{center}
    \begin{tikzpicture}
      \bitcells[A]{0}{00000000000000000000000000000001}{a}
      \bitcells[Bw]{-0.5}{00000000000000000000000000000000}{b}
      \bitcells[C]{-1.0}{00000000000000000000000000000000}{c}
      \bitcells[D]{-1.5}{00000000000000000000000000000000}{d}
      \draw[cellmark] (A-31.south west) rectangle (A-31.north east);
    \end{tikzpicture}
  \end{center}
  \vfill
  \begin{center}\small\textbf{Predict first:} the input has a single 1 bit.\\
    How many 1 bits will the four words have at the end?\end{center}
  \vfill
  \say{Około 6 minut. Najpierw niech zapiszą przewidywanie na kartce. Można liczyć ręcznie, ale Python jest szybszy.}
\end{frame}

\samedot
\begin{frame}[fragile]{Task 4: what to do}
  \begin{columns}[T,onlytextwidth]
    \column{0.62\textwidth}
\begin{lstlisting}[style=py, basicstyle=\color{VSText}\ttfamily\fontsize{6.4}{7.4}\selectfont]
def salsa_quarter_round(a, b, c, d):
    # TODO: the four lines of the Salsa20 quarter round
    #   b ^= rotl(a + d, 7)
    #   c ^= rotl(b + a, 9)
    #   d ^= rotl(c + b, 13)
    #   a ^= rotl(d + c, 18)
    # use add32 for + and rotl32 for rotl;
    # print a, b, c, d after every line
    return a, b, c, d
\end{lstlisting}
    \column{0.35\textwidth}
      \taskhead{Task 4}{Open \texttt{salsa\_qr.py} from \ghlink{Task4}.}
      \stepbox{Ready in the file}{\texttt{add32} and \texttt{rotl32}}
  \end{columns}
  \say{Gotowe add32 i rotl32 są w pliku, więc zadanie nie wymaga zadania 3.}
\end{frame}

\samedot
\begin{frame}{Solution: one quarter round of Salsa20}
  \vfill
  \begin{center}
    % CHECKED: Task4/solution/salsa_qr.py; the same as the example in Bernstein's Salsa20 specification
    \ttfamily\small
    \renewcommand{\arraystretch}{1.35}
    \begin{tabular}{l c c c c}
      \toprule
      & a & b & c & d \\
      \midrule
      start & 0x00000001 & 0x00000000 & 0x00000000 & 0x00000000 \\
      line 1 & 0x00000001 & \cellcolor{AccentOrange!18}0x00000080 & 0x00000000 & 0x00000000 \\
      line 2 & 0x00000001 & 0x00000080 & \cellcolor{AccentOrange!18}0x00010200 & 0x00000000 \\
      line 3 & 0x00000001 & 0x00000080 & 0x00010200 & \cellcolor{AccentOrange!18}0x20500000 \\
      line 4 & \cellcolor{AccentOrange!18}0x08008145 & 0x00000080 & 0x00010200 & 0x20500000 \\
      \bottomrule
    \end{tabular}
  \end{center}
  \vfill
  \begin{center}\small 1 bit in, \textbf{12 bits} out, after one quarter round.\end{center}
  \vfill
  \keymsg{One bit already reaches all four words.}
  \say{Wynik zgadza się z przykładem w specyfikacji Salsa20 Bernsteina (6 + 1 + 2 + 3 = 12 jedynek). Po pełnych 20 rundach różnica obejmie cały blok, zmierzymy to w sekcji 5.}
\end{frame}

% the block function as a flow, top to bottom: #1 = 1 mark the final addition
\newcommand{\blockflow}[1]{%
  \def\flowwide{1}\ifnum#1=1 \def\flowwide{0}\fi
  \begin{tikzpicture}[>=Stealth, node distance=0.42cm]
    \node[cell, font=\small, minimum width=4.6cm] (in) {key + nonce + counter + constants};
    \node[cell, fill=nFill, font=\small, minimum width=4.6cm, below=of in] (st) {initial state (16 words)};
    \node[hwmain, font=\small\bfseries, minimum width=4.6cm, below=of st] (r) {\faCogs\ 20 rounds of quarter rounds};
    \node[circle, draw=nInk, very thick, fill=white, font=\large, inner sep=0pt, minimum size=0.6cm, below=of r] (add) {$+$};
    \node[ksb, font=\small, minimum width=4.6cm, minimum height=0.55cm, below=of add] (out) {64 bytes of keystream};
    \draw[data] (in) -- (st); \draw[data] (st) -- (r); \draw[data] (r) -- (add); \draw[data] (add) -- (out);
    \draw[data] (st.east) -- ++(1.0,0) coordinate (side) |- (add.east);
    \ifnum\flowwide=1
      \node[font=\scriptsize, text=nInk, anchor=west] at ($(side)!0.5!(side |- add.east)$) {add the initial state};
      \draw[data] (out.east) -- ++(0.6,0) node[right, font=\scriptsize, text=nInk] {XOR with 64 bytes of plaintext};
    \fi
    \ifnum#1=1 \draw[annot] ([shift={(-0.1,-0.1)}]add.south west) rectangle ([shift={(0.1,0.1)}]add.north east); \fi
  \end{tikzpicture}}

\begin{frame}{From the state to a block}
  \vfill
  \begin{center}\blockflow{0}\end{center}
  \vfill
  \say{Ten sam przepływ ma ChaCha20. Po co jest dodanie stanu początkowego na końcu, sprawdzimy w sekcji 5 atakiem.}
\end{frame}

\begin{frame}{Salsa20 today}
  \vfill
  \begin{center}
    \begin{tikzpicture}
      \draw[nLine, very thick] (0,0) -- (10,0);
      \foreach \x/\y/\t in {1/2005/{Salsa20 designed}, 5/2008/{eSTREAM recommends Salsa20/12}, 9/today/{NaCl, libsodium (as XSalsa20)}} {
        \fill[nInk] (\x,0) circle (0.1);
        \node[font=\small\bfseries, text=nInk, anchor=south] at (\x,0.2) {\y};
        \node[font=\scriptsize, text=nInk, anchor=north, text width=3.2cm, align=center] at (\x,-0.2) {\t};
      }
    \end{tikzpicture}
  \end{center}
  \vfill
  \small \textbf{eSTREAM}: a European competition for stream ciphers. In 2008 it recommends Salsa20/12 (12 rounds) for software.\\[4pt]
  Used in the NaCl and libsodium libraries (as XSalsa20), for example in Keybase chat.
  \vfill
  \keymsg{ChaCha20 keeps the idea and changes the details.}
  \say{XSalsa20 to Salsa20 z dłuższym nonce. Salsa20 jest nadal bezpieczny, ale nowe systemy wybierają ChaCha20.}
\end{frame}

% ======================================================================
\section[ChaCha20: one block]{ChaCha20: one block}
% ======================================================================

% CHECKED (python3, Task6/solution/block.py and a separate script, RFC 8439):
%   2.1.1 quarter round and every step below; 2.3.2 state, state after 20 rounds, block 10f1e7e4...

\begin{frame}{What ChaCha20 changes}
  \vfill
  \begin{columns}[T,onlytextwidth]
    \column{0.32\textwidth}
      \stepbox{\badge{1}\ A new quarter round}{each word changes twice, rotations 16, 12, 8, 7.}
    \column{0.32\textwidth}
      \stepbox{\badge{2}\ A new layout}{the constants in the first row, the key in rows 2 and 3.}
    \column{0.32\textwidth}
      \stepbox{\badge{3}\ Rounds alternate}{columns, then diagonals.}
  \end{columns}
  \vfill
  \keymsg{More mixing per round, the same three operations.}
  \say{Bernstein pokazał w 2008 roku, że te zmiany szybciej rozprowadzają różnice między bitami, a szyfr nie jest wolniejszy. Każdą z trzech zmian zobaczymy osobno.}
\end{frame}

\newcommand{\chachalayout}{const/c,const/c,const/c,const/c,key/k,key/k,key/k,key/k,key/k,key/k,key/k,key/k,counter/t,nonce/n,nonce/n,nonce/n}
\newcommand{\rfcstate}{61707865/c,3320646e/c,79622d32/c,6b206574/c,03020100/k,07060504/k,0b0a0908/k,0f0e0d0c/k,13121110/k,17161514/k,1b1a1918/k,1f1e1d1c/k,00000001/t,09000000/n,4a000000/n,00000000/n}

\begin{frame}{The ChaCha20 state}
  \vfill
  \begin{center}\statefig[1.9]{\chachalayout}{}\\[6pt] \statelegend\end{center}
  \vfill
  \begin{center}\small constants 128 bits\qquad key 256 bits\qquad counter 32 bits\qquad nonce 96 bits\end{center}
  \vfill
  \say{To wersja z RFC 8439, używana w protokołach internetowych: licznik 32 bity, nonce 96 bitów. Oryginalny ChaCha20 Bernsteina miał po 64 bity, jak Salsa20.}
\end{frame}

\samedot
\begin{frame}{Little-endian}
  \begin{center}\small Four bytes become one word with the \textbf{lowest byte first}:\end{center}
  \vfill
  \begin{center}
    \begin{tikzpicture}[>=Stealth]
      \node[rowlabel, font=\small] at (-0.6,1.2) {bytes of the key};
      \foreach \b [count=\k from 0] in {00,01,02,03} \node[st-k, minimum width=0.9cm, font=\ttfamily\small] (y\k) at ({1.1*\k},1.2) {\b};
      \node[rowlabel, font=\small] at (-0.6,-0.4) {word};
      \node[font=\ttfamily\small, text=nInk, anchor=east] at (-0.18,-0.4) {0x};
      \foreach \b [count=\k from 0] in {03,02,01,00} \node[st-k, minimum width=0.9cm, font=\ttfamily\small] (w\k) at ({1.1*\k+0.3},-0.4) {\b};
      \foreach \k in {0,1,2,3} {\pgfmathtruncatemacro{\j}{3-\k} \draw[data, thick] (y\k.south) -- (w\j.north);}
    \end{tikzpicture}
  \end{center}
  \vfill
  \begin{center}\small In Python:\ \texttt{struct.unpack("<I", b"\textbackslash x00\textbackslash x01\textbackslash x02\textbackslash x03")}\ $\rightarrow$\ \texttt{0x03020100}\\[4pt]
    This is the order used by most processors.\end{center}
  \vfill
  \say{Bez tego slajdu nikt nie odczyta stanu z RFC. < w struct oznacza little-endian, I to słowo 32-bitowe.}
\end{frame}

\samedot
\begin{frame}{The state of the RFC example}
  \begin{center}\small key = \texttt{00 01 02 \ldots\ 1f},\quad counter = 1,\quad nonce = \texttt{00 00 00 09 00 00 00 4a 00 00 00 00}\end{center}
  \vfill
  \begin{center}\statefig[1.9]{\rfcstate}{}\\[6pt] \statelegend\end{center}
  \vfill
  \say{To przykład z RFC 8439, punkt 2.3.2. Te liczby zobaczymy do końca sekcji, a w zadaniu 6 studenci dostaną dokładnie ten wynik.}
\end{frame}

\begin{frame}{The ChaCha20 quarter round}
  \begin{columns}[c,onlytextwidth]
    \column{0.48\textwidth}
      \centering
      {\small quarter round of four words a, b, c, d}\\[4pt]
      \qrcode{0}
    \column{0.49\textwidth}
      \stepbox{Every line}{add, XOR, rotate.\quad \texttt{+=} is addition mod $2^{32}$.}
      \stepbox{Each of the four words changes twice}{In Salsa20 every word changed once.}
  \end{columns}
  \vfill
  \say{Porównać z Salsa20: tam każde słowo zmieniało się raz. Tu a i c dodają, b i d są mieszane XOR i rotacją. Teraz policzymy to na liczbach z RFC.}
\end{frame}

% one step of the quarter round: #1 line, #2 three lines of the calculation, #3 to #6 the words, #7 changed words
\newcommand{\qrstep}[7]{%
  \begin{columns}[c,onlytextwidth]
    \column{0.47\textwidth}
      \centering\qrcode{#1}
    \column{0.5\textwidth}
      \plainbox{\raggedright\ttfamily\small #2}
  \end{columns}
  \vfill
  \begin{center}\abcd{#3}{#4}{#5}{#6}{#7}\end{center}
  \vfill}

\samedot
\begin{frame}{Quarter round: the start}
  \qrstep{0}{RFC 8439, section 2.1.1:\\ the test vector of the quarter round}{0x11111111}{0x01020304}{0x9b8d6f43}{0x01234567}{}
  \say{Wektor testowy z RFC 8439, punkt 2.1.1. Liczymy linijka po linijce; Python wystarcza jako kalkulator.}
\end{frame}

\samedot
\begin{frame}{Quarter round, line 1}
  \qrstep{1}{a = 0x11111111 + 0x01020304 = 0x12131415\\ d = 0x01234567 \textasciicircum{} 0x12131415 = 0x13305172\\ d = rotl(0x13305172, 16)\ \ = 0x51721330}{0x12131415}{0x01020304}{0x9b8d6f43}{0x51721330}{a,d}
  \say{Pierwsza linijka. XOR na 0x12131415 i 0x01234567 widzieliśmy już przy operacji 2. Rotacja o 16 zamienia połówki słowa miejscami.}
\end{frame}

\samedot
\begin{frame}{Quarter round, line 2}
  \qrstep{2}{c = 0x9b8d6f43 + 0x51721330 = 0xecff8273\\ b = 0x01020304 \textasciicircum{} 0xecff8273 = 0xedfd8177\\ b = rotl(0xedfd8177, 12)\ \ = 0xd8177edf}{0x12131415}{0xd8177edf}{0xecff8273}{0x51721330}{b,c}
  \say{Rotacja o 12 bitów to przesunięcie o 3 cyfry hex, więc łatwo ją sprawdzić okiem: edfd8177 staje się d8177edf.}
\end{frame}

\samedot
\begin{frame}{Quarter round, line 3}
  \qrstep{3}{a = 0x12131415 + 0xd8177edf = 0xea2a92f4\\ d = 0x51721330 \textasciicircum{} 0xea2a92f4 = 0xbb5881c4\\ d = rotl(0xbb5881c4, 8)\ \ \ = 0x5881c4bb}{0xea2a92f4}{0xd8177edf}{0xecff8273}{0x5881c4bb}{a,d}
  \say{Rotacja o 8 bitów to przesunięcie o dwie cyfry hex: bb5881c4 staje się 5881c4bb.}
\end{frame}

\samedot
\begin{frame}{Quarter round, line 4}
  \qrstep{4}{c = 0xecff8273 + 0x5881c4bb = 0x4581472e\\ \textnormal{\scriptsize\textcolor{AccentOrange}{(a carry was lost)}}\\ b = 0xd8177edf \textasciicircum{} 0x4581472e = 0x9d9639f1\\ b = rotl(0x9d9639f1, 7)\ \ \ = 0xcb1cf8ce}{0xea2a92f4}{0xcb1cf8ce}{0x4581472e}{0x5881c4bb}{b,c}
  \keymsg{The result matches RFC 8439, section 2.1.1.}
  \say{Przy dodawaniu wypadła jedynka poza 32 bity, to jest właśnie mod 2 do 32. Rotacji o 7 nie da się już sprawdzić okiem, bo 7 nie dzieli się przez 4.}
\end{frame}

\begin{frame}{Task 5: write the quarter round}
  \taskhead{Task 5}{Write the ChaCha20 quarter round on four words. Test it on the example from RFC 8439:}
  \vfill
  \begin{center}
    \ttfamily\small
    \renewcommand{\arraystretch}{1.5}
    \begin{tabular}{l l l l l}
      \normalfont in & 0x11111111 & 0x01020304 & 0x9b8d6f43 & 0x01234567 \\
      \normalfont out & 0xea2a92f4 & 0xcb1cf8ce & 0x4581472e & 0x5881c4bb \\
    \end{tabular}
  \end{center}
  \vfill
  \say{Około 8 minut. Wartości pośrednie są na poprzednich slajdach, przydadzą się, jeśli coś nie wyjdzie.}
\end{frame}

\samedot
\begin{frame}[fragile]{Task 5: what to do}
  \begin{columns}[T,onlytextwidth]
    \column{0.6\textwidth}
\begin{lstlisting}[style=py, basicstyle=\color{VSText}\ttfamily\fontsize{6.4}{7.4}\selectfont]
def quarter_round(a, b, c, d):
    # TODO: the four lines of the ChaCha20 quarter round
    #   a += b; d ^= a; d = rotl(d, 16)
    #   c += d; b ^= c; b = rotl(b, 12)
    #   a += b; d ^= a; d = rotl(d, 8)
    #   c += d; b ^= c; b = rotl(b, 7)
    # use add32 for += and rotl32 for rotl
    return a, b, c, d
\end{lstlisting}
    \column{0.37\textwidth}
      \taskhead{Task 5}{Open \texttt{quarter\_round.py} from \ghlink{Task5}.}
      \stepbox{Ready in the file}{\texttt{add32} and \texttt{rotl32}. The program prints OK or FAIL.}
  \end{columns}
  \say{Częsty błąd: zwykłe + zamiast add32, wtedy liczby rosną ponad 32 bity i test daje FAIL.}
\end{frame}

\samedot
\begin{frame}[fragile]{Solution: write the quarter round}
  \begin{columns}[T,onlytextwidth]
    \column{0.58\textwidth}
% CHECKED: Task5/solution/quarter_round.py prints OK
\begin{lstlisting}[style=py, emph={add32,rotl32}]
def quarter_round(a, b, c, d):
    a = add32(a, b); d ^= a; d = rotl32(d, 16)
    c = add32(c, d); b ^= c; b = rotl32(b, 12)
    a = add32(a, b); d ^= a; d = rotl32(d, 8)
    c = add32(c, d); b ^= c; b = rotl32(b, 7)
    return a, b, c, d
\end{lstlisting}
    \column{0.39\textwidth}
      \stepbox{The output}{\scriptsize\texttt{got:\ \ \ \ \ 0xea2a92f4 0xcb1cf8ce\\ \ \ \ \ \ \ \ \ \ \ 0x4581472e 0x5881c4bb}\\[2pt]
        \texttt{OK}}
  \end{columns}
  \say{Kod sprawdzony. Cztery linijki to cały klocek ChaCha20; reszta to tylko decyzja, na których słowach go wywołać.}
\end{frame}

% the state with the four quarter rounds of one round, one of them lit: #1 = 0 column, 1 diagonal; #2 = lit group (1 to 4)
\newcommand{\roundfig}[2]{%
  \ifnum#1=0
    \def\groups{0/A,1/B,2/C,3/D,4/A,5/B,6/C,7/D,8/A,9/B,10/C,11/D,12/A,13/B,14/C,15/D}
  \else
    \def\groups{0/A,1/B,2/C,3/D,4/D,5/A,6/B,7/C,8/C,9/D,10/A,11/B,12/B,13/C,14/D,15/A}
  \fi
  \begin{tikzpicture}
    \foreach \p/\g [count=\i from 0] in \groups {
      \pgfmathtruncatemacro{\r}{\i/4}\pgfmathtruncatemacro{\c}{mod(\i,4)}
      \node[wcell, minimum width=1.5cm, font=\ttfamily\small] (G\i) at ({1.6*\c},{-0.65*\r}) {\p\ \textbf{\g}};
      \edef\lit{\ifcase#2\or A\or B\or C\or D\fi}
      \ifx\g\lit \draw[cellmark] (G\i.south west) rectangle (G\i.north east); \fi
    }
  \end{tikzpicture}}
\newcommand{\qrcall}[3]{\ifnum#1=#2 \textbf{\textcolor{AccentOrange}{#3}}\else #3\fi}

\begin{frame}{A column round}
  \begin{columns}[c,onlytextwidth]
    \column{0.5\textwidth}
      \centering
      \only<1>{\roundfig{0}{1}}\only<2>{\roundfig{0}{2}}\only<3>{\roundfig{0}{3}}\only<4>{\roundfig{0}{4}}
    \column{0.47\textwidth}
      \ttfamily\small
      \only<1>{\qrcall{1}{1}{A: quarter\_round(0, 4, 8, 12)}}\only<2-4>{A: quarter\_round(0, 4, 8, 12)}\\[3pt]
      \only<2>{\qrcall{1}{1}{B: quarter\_round(1, 5, 9, 13)}}\only<1,3-4>{B: quarter\_round(1, 5, 9, 13)}\\[3pt]
      \only<3>{\qrcall{1}{1}{C: quarter\_round(2, 6, 10, 14)}}\only<1-2,4>{C: quarter\_round(2, 6, 10, 14)}\\[3pt]
      \only<4>{\qrcall{1}{1}{D: quarter\_round(3, 7, 11, 15)}}\only<1-3>{D: quarter\_round(3, 7, 11, 15)}
  \end{columns}
  \vfill
  \defbox{\textbf{Column round}: four quarter rounds, each on one column of the state.}
  \note<4>{\selectlanguage{polish}Cztery quarter round, każda na jednej kolumnie. Po rundzie kolumnowej każde słowo zależy od słów ze swojej kolumny, ale nie od innych kolumn.}
\end{frame}

\samedot
\begin{frame}{A diagonal round}
  \begin{columns}[c,onlytextwidth]
    \column{0.5\textwidth}
      \centering
      \only<1>{\roundfig{1}{1}}\only<2>{\roundfig{1}{2}}\only<3>{\roundfig{1}{3}}\only<4>{\roundfig{1}{4}}
    \column{0.47\textwidth}
      \ttfamily\small
      \only<1>{\qrcall{1}{1}{A: quarter\_round(0, 5, 10, 15)}}\only<2-4>{A: quarter\_round(0, 5, 10, 15)}\\[3pt]
      \only<2>{\qrcall{1}{1}{B: quarter\_round(1, 6, 11, 12)}}\only<1,3-4>{B: quarter\_round(1, 6, 11, 12)}\\[3pt]
      \only<3>{\qrcall{1}{1}{C: quarter\_round(2, 7, 8, 13)}}\only<1-2,4>{C: quarter\_round(2, 7, 8, 13)}\\[3pt]
      \only<4>{\qrcall{1}{1}{D: quarter\_round(3, 4, 9, 14)}}\only<1-3>{D: quarter\_round(3, 4, 9, 14)}
  \end{columns}
  \vfill
  \defbox{\textbf{Diagonal round}: four quarter rounds, each on one diagonal; a diagonal wraps around the edge.}
  \note<4>{\selectlanguage{polish}Przekątna zawija się na brzegu tablicy, jak mod. Każda przekątna ma po jednym słowie z każdej kolumny.}
\end{frame}

\samedot
\begin{frame}{Two rounds: every word meets every other}
  \begin{columns}[c,onlytextwidth]
    \column{0.5\textwidth}
      \centering
      \begin{tikzpicture}
        \foreach \i in {0,...,15} {
          \pgfmathtruncatemacro{\r}{\i/4}\pgfmathtruncatemacro{\c}{mod(\i,4)}
          \node[wcell, minimum width=1.5cm, font=\ttfamily\small] (H\i) at ({1.6*\c},{-0.65*\r}) {\i};
        }
        \draw[cellmark] (H0.south west) rectangle (H0.north east);
        \foreach \i in {4,8,12} \draw[annot, dashed] (H\i.south west) rectangle (H\i.north east);
      \end{tikzpicture}
    \column{0.47\textwidth}
      \stepbox{After a column round}{word 0 depends on words 0, 4, 8, 12}
      \stepbox{After a diagonal round}{word 0 depends on \textbf{all 16 words}}
  \end{columns}
  \vfill
  \defbox{A column round and a diagonal round together are a \textbf{double round}.}
  \keymsg{One double round connects every word with every other word.}
  \say{Słowo 0 po rundzie diagonalnej miesza się ze słowami 5, 10 i 15, a każde z nich wcześniej mieszało się ze swoją kolumną. Dlatego zależy od wszystkich szesnastu.}
\end{frame}

\samedot
\begin{frame}{Twenty rounds}
  \begin{columns}[c,onlytextwidth]
    \column{0.42\textwidth}
      \begin{tikzpicture}
        \codebox{3}{5.4}
        \cl{1}{\kw{for} \_ \kw{in} \fn{range}(10):\ \ \cm{\# 10 double rounds}}
        \cl{2}{\tb column round}
        \cl{3}{\tb diagonal round}
      \end{tikzpicture}
    \column{0.56\textwidth}
      \centering
      {\small The RFC example after 20 rounds:}\\[4pt]
      \statefig[1.55]{837778ab/x,e238d763/x,a67ae21e/x,5950bb2f/x,c4f2d0c7/x,fc62bb2f/x,8fa018fc/x,3f5ec7b7/x,335271c2/x,f29489f3/x,eabda8fc/x,82e46ebd/x,d19c12b4/x,b04e16de/x,9e83d0cb/x,4e3c50a2/x}{}
  \end{columns}
  \vfill
  \begin{center}\small Nothing looks like the constants or the key any more.\end{center}
  \say{Wartości zgadzają się z RFC 8439, punkt 2.3.2. Nie widać już ani stałych, ani klucza z kolejnymi bajtami.}
\end{frame}

\begin{frame}{Add the initial state}
  \begin{center}\small Each word: the word after 20 rounds + the same word of the initial state (mod $2^{32}$).\end{center}
  \vfill
  \begin{center}
    \ttfamily\small
    \renewcommand{\arraystretch}{1.5}
    \begin{tabular}{l r c r c r}
      word 0: & 0x837778ab & + & 0x61707865 & = & 0xe4e7f110 \\
      word 1: & 0xe238d763 & + & 0x3320646e & = & 0x15593bd1 \\
      \ldots & & & & & \\
      word 15: & 0x4e3c50a2 & + & 0x00000000 & = & 0x4e3c50a2 \\
    \end{tabular}
  \end{center}
  \vfill
  \keymsg{Why this addition matters, we will see by attacking it.}
  \say{To krok, który łatwo uznać za zbędny. W sekcji 5 studenci sami zobaczą, co się dzieje bez niego.}
\end{frame}

\samedot
\begin{frame}{From words to bytes}
  \begin{center}\small Each of the 16 words is written as 4 bytes, little-endian:\end{center}
  \vfill
  \begin{center}
    \begin{tikzpicture}[>=Stealth]
      \foreach \w/\bs/\y in {e4e7f110/{10,f1,e7,e4}/0, 15593bd1/{d1,3b,59,15}/-1.0} {
        \node[wcell, minimum width=2.0cm, font=\ttfamily\small] at (0,\y) {0x\w};
        \draw[data] (1.15,\y) -- (2.2,\y);
        \foreach \b [count=\k from 0] in \bs \node[ksb, minimum width=0.6cm, font=\ttfamily\small] at ({2.6+0.68*\k},\y) {\b};
      }
    \end{tikzpicture}
  \end{center}
  \vfill
  \begin{center}\small The block:\ \texttt{10 f1 e7 e4 d1 3b 59 15 50 0f dd 1f a3 20 71 c4 \ldots}\ (64 bytes)\end{center}
  \vfill
  \say{To dokładnie pierwszy wiersz bloku z RFC 8439. Te 64 bajty to blok keystreamu numer 1 dla tego klucza i nonce.}
\end{frame}

\samedot
\begin{frame}{XOR with the plaintext}
  \vfill
  \begin{center}
    % CHECKED: python3, "Ladies and" XOR 10 f1 e7 e4 d1 3b 59 15 50 0f = 5c 90 83 8d b4 48 79 74 3e 6b
    \begin{tikzpicture}
      \node[rowlabel, font=\small] at (-0.5,0.55) {plaintext};
      \foreach \ch [count=\k from 0] in {L,a,d,i,e,s,{\textvisiblespace},a,n,d} \node[font=\ttfamily\small, text=sViolet] at ({0.72*\k},0.55) {\ch};
      \byterow{0}{0}{ptb}{4c,61,64,69,65,73,20,61,6e,64}{0.72}
      \node[rowlabel, font=\small] at (-0.5,-0.7) {keystream};
      \byterow{0}{-0.7}{ksb}{10,f1,e7,e4,d1,3b,59,15,50,0f}{0.72}
      \node[rowlabel, font=\small] at (-0.5,-1.4) {ciphertext};
      \byterow{0}{-1.4}{ctb}{5c,90,83,8d,b4,48,79,74,3e,6b}{0.72}
    \end{tikzpicture}
  \end{center}
  \vfill
  \keymsg{The block function makes the keystream. Encryption is still one XOR.}
  \say{Wszystko poza tym XOR-em to tylko produkcja keystreamu. Wartości policzone moją implementacją, zgodną z RFC.}
\end{frame}

\begin{frame}{Task 6: one full block}
  \taskhead{Task 6}{Write the whole block function of ChaCha20.}
  \vfill
  \begin{columns}[T,onlytextwidth]
    \column{0.24\textwidth}\stepbox{\badge{1}}{the initial state}
    \column{0.24\textwidth}\stepbox{\badge{2}}{10 double rounds}
    \column{0.24\textwidth}\stepbox{\badge{3}}{add the initial state}
    \column{0.24\textwidth}\stepbox{\badge{4}}{64 bytes, little-endian}
  \end{columns}
  \vfill
  \begin{center}\small Test: the example from RFC 8439, the block must start with\ \texttt{\textbf{10 f1 e7 e4}}.\end{center}
  \vfill
  \say{Około 12 minut, najdłuższe zadanie. quarter\_round jest gotowe w pliku.}
\end{frame}

\samedot
\begin{frame}[fragile]{Task 6: what to do}
  \begin{columns}[T,onlytextwidth]
    \column{0.64\textwidth}
\begin{lstlisting}[style=py, basicstyle=\color{VSText}\ttfamily\fontsize{6.2}{7.1}\selectfont]
def chacha20_block(key, counter, nonce):
    # TODO 1: the initial state, 16 words:
    #         CONSTANTS + 8 key words + counter + 3 nonce words
    #         struct.unpack("<8I", key) reads 8 little-endian words
    start = []

    # TODO 2: copy the state and run 10 double rounds on the copy:
    #   a column round:
    #     (0, 4, 8, 12) (1, 5, 9, 13) (2, 6, 10, 14) (3, 7, 11, 15)
    #   a diagonal round:
    #     (0, 5, 10, 15) (1, 6, 11, 12) (2, 7, 8, 13) (3, 4, 9, 14)
    s = start[:]

    # TODO 3: add the initial state to the result, word by word, with add32

    # TODO 4: return the 16 words as 64 bytes: struct.pack("<16I", *s)
    return b""
\end{lstlisting}
    \column{0.33\textwidth}
      \taskhead{Task 6}{Open \texttt{block.py} from \ghlink{Task6}.}
      \stepbox{Ready in the file}{\texttt{quarter\_round(s, a, b, c, d)}: it changes the list \texttt{s} in place.}
  \end{columns}
  \say{Uwaga: quarter\_round w tym pliku przyjmuje listę i cztery indeksy, a nie cztery wartości jak w zadaniu 5. Tak jest wygodniej dla całego stanu. Indeksy kolumn i przekątnych są w komentarzu w pliku.}
\end{frame}

\samedot
\begin{frame}[fragile]{Solution: one full block}
  \begin{columns}[T,onlytextwidth]
    \column{0.64\textwidth}
% CHECKED: Task6/solution/block.py prints OK; the slide shows 2 of the 8 calls of quarter_round
\begin{lstlisting}[style=py, emph={quarter_round,add32,list,range}, basicstyle=\color{VSText}\ttfamily\fontsize{6.4}{7.4}\selectfont]
start = CONSTANTS + list(struct.unpack("<8I", key)) \
        + [counter] + list(struct.unpack("<3I", nonce))
s = start[:]
for _ in range(10):
    quarter_round(s, 0, 4, 8, 12)    # column round
    ...                               # 3 more columns
    quarter_round(s, 0, 5, 10, 15)   # diagonal round
    ...                               # 3 more diagonals
for i in range(16):
    s[i] = add32(s[i], start[i])
return struct.pack("<16I", *s)
\end{lstlisting}
    \column{0.33\textwidth}
      \stepbox{\badge{1} to \badge{4}}{the four parts from the task, in order.}
      \stepbox{Shortened on the slide}{the solution file has all 8 calls of \texttt{quarter\_round}.}
  \end{columns}
  \say{Kod sprawdzony. Na slajdzie skrócone: w pliku rozwiązania jest wszystkie osiem wywołań quarter\_round.}
\end{frame}

\samedot
\begin{frame}{Solution: the output}
  \vfill
  \begin{center}
    \ttfamily\small
    \begin{tabular}{l l}
      got: & 10f1e7e4d13b5915500fdd1fa32071c4c7d1f4c733c068030422aa9ac3d46c4e\ldots \\
      expected: & 10f1e7e4d13b5915500fdd1fa32071c4c7d1f4c733c068030422aa9ac3d46c4e\ldots \\
      & \ok\ \normalfont OK \\
    \end{tabular}
  \end{center}
  \vfill
  \keymsg{Your ChaCha20 gives the same block as the standard.}
  \say{Wynik ma 64 bajty, na slajdzie pierwsze 32. Implementacja studentów zgadza się z każdą biblioteką kryptograficzną na świecie.}
\end{frame}

\begin{frame}[fragile]{All of ChaCha20 on one slide}
  \begin{columns}[T,onlytextwidth]
    \column{0.5\textwidth}
\begin{lstlisting}[style=py, emph={add32,rotl32,list,range}, basicstyle=\color{VSText}\ttfamily\fontsize{5.6}{6.6}\selectfont]
def quarter_round(s, a, b, c, d):
    s[a] = add32(s[a], s[b]); s[d] ^= s[a]; s[d] = rotl32(s[d], 16)
    s[c] = add32(s[c], s[d]); s[b] ^= s[c]; s[b] = rotl32(s[b], 12)
    s[a] = add32(s[a], s[b]); s[d] ^= s[a]; s[d] = rotl32(s[d], 8)
    s[c] = add32(s[c], s[d]); s[b] ^= s[c]; s[b] = rotl32(s[b], 7)


def initial_state(key, counter, nonce):
    return ([0x61707865, 0x3320646E, 0x79622D32, 0x6B206574]
            + list(struct.unpack("<8I", key))
            + [counter]
            + list(struct.unpack("<3I", nonce)))
\end{lstlisting}
    \column{0.48\textwidth}
\begin{lstlisting}[style=py, emph={add32,initial_state,run_rounds,range,len,bytes,block,bytearray}, basicstyle=\color{VSText}\ttfamily\fontsize{5.6}{6.6}\selectfont]
def block(key, counter, nonce, rounds=20):
    start = initial_state(key, counter, nonce)
    s = start[:]
    run_rounds(s, rounds)
    for i in range(16):
        s[i] = add32(s[i], start[i])
    return struct.pack("<16I", *s)

def chacha20(key, nonce, counter, data):
    out = bytearray()
    for i in range(len(data)):
        if i % 64 == 0:
            ks = block(key, counter + i // 64, nonce)
        out.append(data[i] ^ ks[i % 64])
    return bytes(out)
\end{lstlisting}
      {\scriptsize\textcolor{gray}{Not shown: \texttt{add32}, \texttt{rotl32} and \texttt{run\_rounds} (the column and diagonal rounds).}}
  \end{columns}
  \vfill
  \keymsg{This is \texttt{chacha.py} from Task 1.}
  \say{Każdą linijkę już widzieliśmy. To jest prawdziwy kod z chacha.py z zadania 1, skrócony tylko o add32, rotl32 i run\_rounds; można go teraz otworzyć.}
\end{frame}

% ======================================================================
\section[Why twenty rounds?]{Why twenty rounds?}
% ======================================================================

\begin{frame}{How many rounds are enough?}
  \vfill
  \begin{columns}[c,onlytextwidth]
    \column{0.3\textwidth}
      \centering{\fontsize{70}{70}\selectfont\textcolor{nLine}{?\,?}}
    \column{0.66\textwidth}
      \Large Why 20 rounds, and not 2?\\[0.3cm]
      Why add the initial state at the end?\\[0.3cm]
      \normalsize Let us \textbf{measure} first, then \textbf{attack}.
  \end{columns}
  \vfill
  \say{Pytanie do sali przed zadaniem: ile rund by wystarczyło i jak to sprawdzić?}
\end{frame}

\begin{frame}{Task 7: watch one bit spread}
  \taskhead{Task 7}{Two blocks with the same key and nonce. The counters differ in \textbf{one bit}.
    How many of the 512 output bits are different?}
  \vfill
  \begin{columns}[c,onlytextwidth]
    \column{0.55\textwidth}
      \centering
      \begin{tikzpicture}
        \foreach \row/\y in {0/0, 1/-0.75} {
          \foreach \k in {0,...,31} \node[cell, minimum width=0.17cm, minimum height=0.5cm, inner sep=0pt] at ({0.19*\k},\y) {};
        }
        \node[rowlabel, font=\scriptsize] at (-0.2,0) {block 1};
        \node[rowlabel, font=\scriptsize] at (-0.2,-0.75) {block 2};
        \foreach \k in {3,9,10,17,22,28,30} \fill[AccentOrange!60] ({0.19*\k-0.085},-1.0) rectangle ({0.19*\k+0.085},-0.5);
        \node[font=\scriptsize, text=AccentOrange, anchor=north] at (3.0,-1.1) {different bits: how many?};
      \end{tikzpicture}
    \column{0.42\textwidth}
      \small
      \begin{tabular}{ll}
        after 1 round: & \textbf{?} \\
        after 2 rounds: & \textbf{?} \\
        after 20 rounds: & \textbf{?} \\
      \end{tabular}\\[6pt]
      \textbf{Predict first}, on paper. Then measure.
  \end{columns}
  \vfill
  \say{Około 6 minut. Zebrać przewidywania od kilku osób i zapisać na tablicy, zanim uruchomią program. (Pomarańczowe kratki na rysunku to tylko ilustracja.)}
\end{frame}

\samedot
\begin{frame}[fragile]{Task 7: what to do}
  \begin{columns}[T,onlytextwidth]
    \column{0.66\textwidth}
\begin{lstlisting}[style=py, emph={block,range,print}, basicstyle=\color{VSText}\ttfamily\fontsize{5.9}{6.9}\selectfont]
for rounds in [1, 2, 3, 4, 5, 6, 8, 12, 20]:
    total = 0
    for _ in range(TRIALS):
        key = random.randbytes(32)
        nonce = random.randbytes(12)
        counter = random.getrandbits(32)
        out1 = block(key, counter, nonce, rounds)
        out2 = block(key, counter ^ 1, nonce, rounds)
        # TODO: count how many of the 512 bits differ
        #       between out1 and out2, add it to total
    print(f"{rounds:2} rounds: {total / TRIALS:6.1f} of 512 bits changed")
\end{lstlisting}
    \column{0.31\textwidth}
      \taskhead{Task 7}{Open \texttt{avalanche.py} from \ghlink{Task7}.}
      \stepbox{Hints}{\texttt{out1[i] \textasciicircum{} out2[i]} has a 1 where the two bytes differ.\\[2pt]
        \texttt{(x >{}> k) \& 1} is bit \texttt{k} of \texttt{x}.}
  \end{columns}
  \say{chacha.py w folderze ma parametr rounds, więc można wybrać dowolną liczbę rund. Program liczy krócej niż sekundę. (x >{}> k) \& 1: przesuwamy x o k bitów w prawo i bierzemy najmłodszy bit, czyli bit numer k.}
\end{frame}

\samedot
\begin{frame}[fragile]{Solution: watch one bit spread}
% CHECKED: Task7/solution/avalanche.py, random.seed(1), 200 trials:
%          1: 26.7, 2: 169.7, 3: 256.3, 4: 255.2, 5: 253.6, 6: 256.1, 8: 255.3, 12: 256.1, 20: 256.3
\begin{lstlisting}[style=py, emph={range}]
for i in range(64):
    x = out1[i] ^ out2[i]
    for k in range(8):
        total += (x >> k) & 1
\end{lstlisting}
  \vfill
  \begin{center}
    \begin{tikzpicture}
      \foreach \r/\v [count=\k from 0] in {1/26.7, 2/169.7, 3/256.3, 4/255.2, 8/255.3, 20/256.3}
        \hbarrow{-0.42*\k}{\r\ \ifnum\r=1 round\else rounds\fi}{\v}{0.022}{\v}
      \draw[DeepBlue, very thick, dashed] ({256*0.022},0.3) -- ({256*0.022},-2.4) node[below, font=\scriptsize] {256 = half of 512};
    \end{tikzpicture}\\[2pt]
    {\scriptsize\textcolor{gray}{average over 200 random keys, of 512 bits}}
  \end{center}
  \vfill
  \say{Kod sprawdzony, random.seed(1), 200 prób. Pełne wyniki z programu: 5 rund 253,6, 6 rund 256,1, 12 rund 256,1. Porównać z przewidywaniami z tablicy.}
\end{frame}

\begin{frame}{Why exactly half?}
  \begin{columns}[c,onlytextwidth]
    \column{0.6\textwidth}
      \small Two random blocks of 512 bits:\\[3pt]
      each bit is the same with chance 1/2, different with chance 1/2\\[3pt]
      so on average $512 / 2 = \mathbf{256}$ bits differ.
    \column{0.36\textwidth}
      \centering
      \begin{tikzpicture}
        \node[circle, draw=nInk, very thick, fill=nFill, minimum size=1.4cm, font=\large\bfseries, text=nInk] at (0,0) {0};
        \node[circle, draw=nInk, very thick, fill=nFill, minimum size=1.4cm, font=\large\bfseries, text=nInk] at (1.8,0) {1};
        \node[font=\scriptsize, text=nInk] at (0.9,-1.05) {a coin for every bit};
      \end{tikzpicture}
  \end{columns}
  \vfill
  \defbox{\textbf{Avalanche}: one changed input bit changes about half of the output bits.}
  \begin{center}\small After 3 rounds, one bit of difference looks like two unrelated blocks.\end{center}
  \keymsg{One changed input bit changes about half of the output bits.}
  \say{Taki efekt nazywa się efektem lawiny: mała zmiana na wejściu zmienia połowę wyjścia. Więcej niż połowa też byłaby zła, bo wtedy wynik byłby przewidywalny w drugą stronę.}
\end{frame}

\begin{frame}{Looks random is not the same as secure}
  \vfill
  \begin{center}
    \begin{tikzpicture}
      \draw[nLine, very thick] (0,0) -- (11.4,0);
      \foreach \r in {1,...,20} {
        \draw[nLine, thick] ({0.57*\r},-0.08) -- ({0.57*\r},0.08);
        \node[font=\tiny, text=nInk, anchor=north] at ({0.57*\r},-0.1) {\r};
      }
      \node[font=\scriptsize, text=nInk, anchor=north] at (5.7,-0.45) {rounds};
      \fill[nLine] (0.57,0.35) rectangle ({0.57*3},0.65);
      \node[font=\scriptsize, text=nInk, anchor=west] at ({0.57*3+0.1},0.5) {our measurement: after 3 rounds we see nothing};
      \fill[nInk] (0.57,0.95) rectangle ({0.57*7},1.25);
      \node[font=\scriptsize, text=nInk, anchor=west] at ({0.57*7+0.1},1.1) {researchers: attacks on 6 and 7 rounds (2008)};
      \draw[annot] ({0.57*20-0.2},-0.3) rectangle ({0.57*20+0.2},0.3);
      \node[font=\scriptsize\bfseries, text=AccentOrange, anchor=south] at ({0.57*20},0.35) {ChaCha20};
    \end{tikzpicture}
  \end{center}
  \vfill
  \begin{center}\small A simple count is a lower limit. Clever attacks see much further.\end{center}
  \vfill
  \keymsg{Passing one test does not prove a cipher secure.}
  \say{Aumasson, Fischer, Khazaei, Meier i Rechberger, FSE 2008: ataki na 6 i 7 rund ChaCha. Ich metody szukają różnic dużo sprytniej niż liczenie bitów.}
\end{frame}

\samedot
\begin{frame}{Differential cryptanalysis: the idea}
  \begin{columns}[c,onlytextwidth]
    \column{0.55\textwidth}
      \footnotesize
      \begin{enumerate}\setlength{\itemsep}{4pt}
        \item Choose two inputs with a chosen difference, for example one bit of the counter.
        \item Follow the difference through the rounds.
        \item Find an output bit that changes more often or less often than half the time.
      \end{enumerate}
    \column{0.42\textwidth}
      \centering
      \begin{tikzpicture}
        \foreach \r/\n in {0/1, 1/3, 2/9} {
          \foreach \k in {0,...,11} \node[cell, minimum width=0.24cm, minimum height=0.3cm, inner sep=0pt] (d\r-\k) at ({0.27*\k},{-0.7*\r}) {};
          \node[font=\tiny, text=gray, anchor=east] at (-0.2,{-0.7*\r}) {round \r};
        }
        \fill[AccentOrange!60] (d0-5.south west) rectangle (d0-5.north east);
        \foreach \k in {4,5,8} \fill[AccentOrange!60] (d1-\k.south west) rectangle (d1-\k.north east);
        \foreach \k in {0,1,3,4,6,8,9,10,11} \fill[AccentOrange!60] (d2-\k.south west) rectangle (d2-\k.north east);
        \draw[annot] (d2-2.south west) rectangle (d2-2.north east);
        \node[font=\tiny, text=AccentOrange, anchor=north] at (d2-2.south) {predictable};
      \end{tikzpicture}
  \end{columns}
  \vfill
  \defbox{\textbf{Differential cryptanalysis}: if such a bit exists, the output is not random, and the attacker learns about the key.}
  \say{Nasze zadanie 7 to bardzo prosta wersja tej metody: patrzyliśmy na sumę bitów, a nie na pojedyncze bity. Ataki z 2008 roku szukają pojedynczych przewidywalnych bitów po kilku rundach. (Rysunek jest ilustracją, nie pomiarem.)}
\end{frame}

\samedot
\begin{frame}{Linear cryptanalysis: the idea}
  \vfill
  \defbox{\textbf{Linear cryptanalysis}: look for a XOR of a few input bits and output bits that is 0 more often than half the time.}
  \vfill
  \begin{columns}[T,onlytextwidth]
    \column{0.48\textwidth}
      \stepbox{XOR and rotation}{keep such relations exact}
    \column{0.48\textwidth}
      \stepbox{Addition}{breaks them, a little in every round}
  \end{columns}
  \vfill
  \begin{center}\small After enough rounds, no useful relation is left.\end{center}
  \vfill
  \say{To tutaj przydaje się dodawanie z sekcji 2: bez niego cały szyfr byłby układem równań XOR i dałby się rozwiązać.}
\end{frame}

\begin{frame}{The best known attacks}
  \vfill
  \begin{center}
    % source: Aumasson, Fischer, Khazaei, Meier, Rechberger, "New Features of Latin Dances", FSE 2008
    \small
    \renewcommand{\arraystretch}{1.5}
    \begin{tabular}{c l l}
      \toprule
      rounds & best known attack & brute force \\
      \midrule
      6 & about $2^{139}$ operations (2008) & $2^{256}$ \\
      7 & about $2^{248}$ operations (2008) & $2^{256}$ \\
      8 & no attack known & $2^{256}$ \\
      \rowcolor{AccentOrange!12} 20 & no attack known & $2^{256}$ \\
      \bottomrule
    \end{tabular}
  \end{center}
  \vfill
  \defbox{\textbf{Brute force}: try every possible key, $2^{256}$ of them.}
  \keymsg{Even 7 rounds are far from broken in practice. 20 rounds is a safety margin.}
  \say{2 do 139 operacji to nadal nieosiągalne, ale to jest atak: lepszy niż zgadywanie. Późniejsze prace poprawiły te liczby dla 7 rund, ale 8 rund nadal nie jest złamane.}
\end{frame}

\samedot
\begin{frame}{ChaCha8, ChaCha12, ChaCha20}
  \vfill
  \begin{center}
    % CHECKED (draft, author's laptop, i5-8265U): the implementation from the tasks, only the keystream, 1 MiB
    \small
    \renewcommand{\arraystretch}{1.5}
    \begin{tabular}{l c c r}
      \toprule
      variant & rounds & rounds above the best attack & 1 MiB in plain Python \\
      \midrule
      ChaCha8 & 8 & 1 & 1.00 s \\
      ChaCha12 & 12 & 5 & 1.44 s \\
      \rowcolor{AccentOrange!12} ChaCha20 & 20 & 13 & 2.24 s \\
      \bottomrule
    \end{tabular}
  \end{center}
  \vfill
  \keymsg{Fewer rounds: faster. More rounds: a bigger safety margin.}
  \say{Czas zmierzony na laptopie z i5-8265U tą samą implementacją co w zadaniach, tylko generowanie keystreamu. W C proporcje są podobne, bo rundy to prawie cały koszt. ChaCha20 wyróżniony jako standard.}
\end{frame}

\samedot
\begin{frame}{Fewer rounds in real systems}
  \vfill
  \begin{center}
    \small
    \renewcommand{\arraystretch}{1.8}
    \begin{tabular}{c l l}
      \textcolor{nInk}{\Large\faLock} & \textbf{ChaCha20} & the standard choice on the internet \\
      \textcolor{nInk}{\Large\faMobile} & \textbf{ChaCha12} & Adiantum: disk encryption on cheap Android phones (2019) \\
      \textcolor{nInk}{\Large\faFile} & \textbf{ChaCha8} & proposed in ``Too Much Crypto'' (Aumasson, 2019), not used by the big standards \\
    \end{tabular}
  \end{center}
  \vfill
  \say{Adiantum Google opracował dla tanich telefonów, na których zwykłe szyfrowanie dysku było za wolne. Wybrał 12 rund, bo dysk musi być szybki. Protokoły sieciowe zostały przy 20.}
\end{frame}

\begin{frame}{Rounds can be undone}
  \begin{center}\small Every step of a round has an inverse:\end{center}
  \vfill
  \begin{center}
    \ttfamily\small
    \renewcommand{\arraystretch}{1.6}
    \begin{tabular}{l c l}
      a = a + b & $\longleftrightarrow$ & a = a - b \\
      d = d \textasciicircum{} a & $\longleftrightarrow$ & d = d \textasciicircum{} a \\
      d = rotl(d, 16) & $\longleftrightarrow$ & d = rotr(d, 16) \\
    \end{tabular}
  \end{center}
  \vfill
  \begin{center}\small Whoever knows the state after 20 rounds can go back to the start.\end{center}
  \vfill
  \keymsg{The rounds hide nothing on their own.}
  \say{Odwracalność jest potrzebna, żeby stan nie tracił informacji. Ale to znaczy też, że same rundy nie ukrywają klucza. Pytanie do sali: co by się stało bez dodania stanu początkowego?}
\end{frame}

\begin{frame}{Task 8: why add the initial state?}
  \taskhead{Task 8}{A \textbf{broken} ChaCha20: 20 rounds, but no addition at the end. You intercepted one output block.}
  \vfill
  \begin{center}
    \statefig[1.6]{ed7caf84/x,ec59ab67/x,182353f4/x,51b53075/x,afb4b0a5/x,9d9a8a93/x,dd172de6/x,23a7a16e/x,38e2e4ca/x,a52d3212/x,2d79c87f/x,382257ec/x,83607414/x,11654b3c/x,a14846fd/x,36134694/x}{}
  \end{center}
  \vfill
  \begin{center}\Large\textbf{Find the secret key.}\end{center}
  \vfill
  \say{Około 10 minut. Napastnik zna tylko ten blok. Nie zna klucza, a mimo to go odzyska.}
\end{frame}

\samedot
\begin{frame}[fragile]{Task 8: what to do}
  \begin{columns}[T,onlytextwidth]
    \column{0.66\textwidth}
\begin{lstlisting}[style=py, emph={inverse_quarter_round,range}, basicstyle=\color{VSText}\ttfamily\fontsize{5.9}{6.9}\selectfont]
def inverse_quarter_round(s, a, b, c, d):
    # TODO: undo the quarter round. Go through its 12 steps backwards:
    #   a += b; d ^= a; d = rotl(d, 16)
    #   c += d; b ^= c; b = rotl(b, 12)
    #   a += b; d ^= a; d = rotl(d, 8)
    #   c += d; b ^= c; b = rotl(b, 7)
    # The last step goes first: b = rotr32(b, 7), then b ^= s[c],
    # then c = sub32(c, d), ...
    # Use s[a], s[b], s[c], s[d] (the list s is changed in place).
    pass


s = output[:]
for _ in range(10):
    # undo one diagonal round, then one column round
    inverse_quarter_round(s, 3, 4, 9, 14)
    ...                                    # 7 more calls, ready
\end{lstlisting}
    \column{0.31\textwidth}
      \taskhead{Task 8}{Open \texttt{invert.py} from \ghlink{Task8}.}
      \stepbox{Ready in the file}{\texttt{sub32} and \texttt{rotr32}, and the loop that undoes the rounds.}
      \footnotesize The program prints rows 1 to 3 of the state.
  \end{columns}
  \say{Kolejność jest odwrócona dwa razy: linijki od ostatniej, a w każdej linijce kroki od ostatniego. Pętla cofająca rundy jest gotowa.}
\end{frame}

\samedot
\begin{frame}[fragile]{Solution: the inverse quarter round}
  \begin{columns}[T,onlytextwidth]
    \column{0.64\textwidth}
% CHECKED: Task8/solution/invert.py prints the constants and the key below
\begin{lstlisting}[style=py, emph={rotr32,sub32}, basicstyle=\color{VSText}\ttfamily\fontsize{6.2}{7.2}\selectfont]
def inverse_quarter_round(s, a, b, c, d):
    s[b] = rotr32(s[b], 7);  s[b] ^= s[c]; s[c] = sub32(s[c], s[d])
    s[d] = rotr32(s[d], 8);  s[d] ^= s[a]; s[a] = sub32(s[a], s[b])
    s[b] = rotr32(s[b], 12); s[b] ^= s[c]; s[c] = sub32(s[c], s[d])
    s[d] = rotr32(s[d], 16); s[d] ^= s[a]; s[a] = sub32(s[a], s[b])
\end{lstlisting}
    \column{0.33\textwidth}
      \centering
      {\scriptsize the original, read from the bottom:}\\[3pt]
      \begin{tikzpicture}[scale=0.8, every node/.append style={transform shape}]
        \codebox{4}{5.0}
        \cl{1}{a += b;\ \ d \textasciicircum= a;\ \ d = rotl(d, 16)}
        \cl{2}{c += d;\ \ b \textasciicircum= c;\ \ b = rotl(b, 12)}
        \cl{3}{a += b;\ \ d \textasciicircum= a;\ \ d = rotl(d, 8)}
        \cl{4}{c += d;\ \ b \textasciicircum= c;\ \ b = rotl(b, 7)}
        \draw[->, very thick, AccentOrange] (5.8,-1.3) -- (5.8,-0.2);
      \end{tikzpicture}
  \end{columns}
  \vfill
  \say{Porównać z quarter round: każda linijka od końca. Kod sprawdzony.}
\end{frame}

\samedot
\begin{frame}{Solution: the key is back}
  \vfill
  \begin{center}
    \small
    \renewcommand{\arraystretch}{1.6}
    \begin{tabular}{l l l}
      constants: & \texttt{61707865 3320646e 79622d32 6b206574} & \textcolor{gray}{"expand 32-byte k"} \\
      key: & \tikz[baseline=-0.6ex]\node[actorlight=sAmber, font=\ttfamily\small, inner sep=3pt] {b'correct horse battery staple!!!!'}; & \\
    \end{tabular}
  \end{center}
  \vfill
  \keymsg{Without the final addition, one block gives away the key.}
  \say{Stałe na początku to potwierdzenie, że cofnęliśmy rundy poprawnie. Kod sprawdzony.}
\end{frame}

\begin{frame}{Why the addition stops the attack}
  \begin{columns}[c,onlytextwidth]
    \column{0.5\textwidth}
      \centering\resizebox{!}{0.6\textheight}{\blockflow{1}}
    \column{0.47\textwidth}
      \begin{center}\large output = rounds(start) + start\end{center}
      \small To go back, the attacker must remove \texttt{start} first.\\[3pt]
      But \texttt{start} contains the key: the attacker does not know it.\\[3pt]
      The rounds can be undone. Undoing rounds plus addition needs the key.
      \vspace{4pt}
      \defbox{A \textbf{one-way} function: easy to compute forward, impossible in practice to undo.}
  \end{columns}
  \vfill
  \keymsg{Adding the initial state turns a reversible mixer into a one-way function.}
  \say{Funkcja jednokierunkowa: łatwo liczyć w przód, nie da się praktycznie cofnąć. Ten sam pomysł jest w wielu funkcjach skrótu.}
\end{frame}

% ======================================================================
\section[Fast on any processor]{Fast on any processor}
% ======================================================================

% a processor chip: \chip{position}{1 = with an AES unit}
\newcommand{\chip}[2]{%
  \begin{scope}[shift={(#1)}]
    \draw[nInk, fill=nInk, rounded corners=2pt] (0,0) rectangle (2.4,1.9);
    \foreach \y in {0.25,0.55,...,1.7} {\draw[nLine, thick] (-0.18,\y) -- (0,\y); \draw[nLine, thick] (2.4,\y) -- (2.58,\y);}
    \node[draw=white, fill=nInk!70, rounded corners=1pt, font=\tiny, text=white, minimum width=0.9cm] at (0.6,1.2) {CPU};
    \ifnum#2=1
      \node[draw=white, fill=nInk!70, rounded corners=1pt, font=\tiny\bfseries, text=white, minimum width=0.9cm] (aes) at (1.75,1.2) {AES};
      \draw[annot] ([shift={(-0.07,-0.07)}]aes.south west) rectangle ([shift={(0.07,0.07)}]aes.north east);
    \else
      \node[draw=white, dashed, rounded corners=1pt, font=\tiny, text=white!60, minimum width=0.9cm] at (1.75,1.2) {no AES};
    \fi
    \node[draw=white, fill=nInk!70, rounded corners=1pt, font=\tiny, text=white, minimum width=2.0cm] at (1.2,0.45) {memory};
  \end{scope}}

\begin{frame}{Many processors have no AES instructions}
  \defbox{\textbf{AES}: the most used \textbf{block cipher} today (it encrypts blocks of 16 bytes).}
  \vfill
  \begin{center}
    \begin{tikzpicture}
      \chip{0,0}{1}
      \node[font=\small, text=nInk, anchor=north] at (1.2,-0.15) {\faLaptop\ laptops, new phones};
      \chip{6,0}{0}
      \node[font=\small, text=nInk, anchor=north] at (7.2,-0.15) {\faMobile\ cheap phones, watches, routers};
    \end{tikzpicture}
  \end{center}
  \vfill
  \begin{center}\small Laptops and new phones have special AES instructions in the processor.\\
    Cheap phones, watches, routers and small devices often do not.\end{center}
  \vfill
  \say{AES to szyfr blokowy, który szyfruje bloki po 16 bajtów. Instrukcje AES na x86 nazywają się AES-NI. Na ARM są rozszerzeniem kryptograficznym, które tanie procesory często pomijają.}
\end{frame}

\samedot
\begin{frame}{AES in software: tables and timing}
  \begin{columns}[c,onlytextwidth]
    \column{0.42\textwidth}
      \centering
      \begin{tikzpicture}[>=Stealth]
        \foreach \k in {0,...,7} \node[cell, minimum width=1.6cm, minimum height=0.34cm] (t\k) at (0,{-0.36*\k}) {};
        \node[font=\scriptsize\bfseries, text=nInk, anchor=south] at (0,0.2) {a table in memory};
        \draw[cellmark] (t5.south west) rectangle (t5.north east);
        \node[actorlight=sAmber, font=\scriptsize\bfseries] (k) at (-2.2,{-0.36*5}) {\faKey\ key};
        \draw[data] (k) -- (t5.west);
        \node[text=nInk, font=\huge] at (2.0,-1.3) {\faStopwatch};
      \end{tikzpicture}
    \column{0.55\textwidth}
      \small
      Fast AES without special instructions uses \textbf{tables in memory}.\\[3pt]
      Which table entry is read depends on the key.\\[3pt]
      Reading memory takes more or less time, depending on the cache.\\[3pt]
      An attacker who measures the time can learn about the key.
  \end{columns}
  \vfill
  \defbox{\textbf{Timing attack}: the attacker learns a secret by measuring how long the computation takes.}
  \keymsg{Fast software AES can leak the key through timing.}
  \say{Cache to mała, szybka pamięć w procesorze: dane, które w niej są, czyta się szybciej. To atak czasowy: napastnik nie łamie matematyki, mierzy czas. Bezpieczne wersje AES w programie istnieją, ale są dużo wolniejsze.}
\end{frame}

\begin{frame}{ChaCha20 needs nothing special}
  \vfill
  \begin{center}
    \small
    \renewcommand{\arraystretch}{1.9}
    \begin{tabular}{c l l}
      \textcolor{nInk}{\Large\faCalculator} & \textbf{only add, rotate, XOR} & every processor has them \\
      \textcolor{nInk}{\Large\faStopwatch} & \textbf{no tables in memory} & the time does not depend on the key \\
      \textcolor{nInk}{\Large\faTh} & \textbf{the same work on many words} & fits the vector instructions of phones \\
    \end{tabular}
  \end{center}
  \vfill
  \keymsg{ChaCha20 is fast and safe in plain software.}
  \say{Instrukcje wektorowe wykonują to samo działanie na kilku słowach naraz, a cztery quarter round w jednej rundzie są od siebie niezależne. Ma je prawie każdy procesor w telefonie.}
\end{frame}

\samedot
\begin{frame}{Measured on a laptop}
  \vfill
  \begin{columns}[T,onlytextwidth]
    \column{0.5\textwidth}
      \centering{\small\bfseries with AES instructions}\\[4pt]
      \begin{tikzpicture}
        \hbarrow{0.9}{AES-128-CTR}{5.7}{0.55}{5.7 GB/s}
        \hbarrow{0.3}{ChaCha20}{3.0}{0.55}{3.0 GB/s}
        \hbarrow{-0.3}{RC4}{0.57}{0.55}{0.57 GB/s}
      \end{tikzpicture}
    \column{0.48\textwidth}
      \centering{\small\bfseries without them}\\[4pt]
      \begin{tikzpicture}
        \hbarrow{0.9}{AES-128-CTR}{0.58}{0.55}{0.58 GB/s}
        \hbarrow{0.3}{ChaCha20}{1.5}{0.55}{1.5 GB/s}
        \draw[annot] (-1.75,0.06) rectangle ({1.5*0.55+1.35},0.54);
      \end{tikzpicture}
  \end{columns}
  \vfill
  \begin{center}\scriptsize\textcolor{gray}{OpenSSL 3.0.13, Intel Core i5-8265U, AES instructions switched off by \texttt{OPENSSL\_ia32cap}}\end{center}
  \keymsg{Without AES instructions, ChaCha20 is about 2.6 times faster than AES.}
  % CHECKED (draft): openssl speed -evp on the author's laptop; 1.5 / 0.58 = 2.6
  \say{Uczciwie: z instrukcjami AES jest prawie 2 razy szybszy od ChaCha20. Bez nich to ChaCha20 wygrywa. Wynik ChaCha20 też spadł po wyłączeniu, bo ta sama maska wyłącza więcej niż tylko AES; nie wpływa to na wniosek.}
\end{frame}

\begin{frame}{Faster than RC4, without its weaknesses}
  \vfill
  \begin{center}
    \small
    \renewcommand{\arraystretch}{1.6}
    \begin{tabular}{l l l}
      \toprule
      & \textbf{RC4} & \textbf{ChaCha20} \\
      \midrule
      speed (laptop) & 0.57 GB/s & 3.0 GB/s \\
      output bias & \nok\ yes, measurable & \ok\ none known \\
      nonce & \nok\ none, glued on by hand & \ok\ built in \\
      random access & \nok\ no & \ok\ yes, by counter \\
      \bottomrule
    \end{tabular}
  \end{center}
  \vfill
  \keymsg{ChaCha20 does what RC4 did, faster and safely.}
  \say{Obciążenie (bias) to odchylenie wyjścia od losowości: w RC4 drugi bajt keystreamu jest zerem dwa razy częściej, niż powinien. W ChaCha20 nie jest znane żadne takie odchylenie dla pełnych 20 rund.}
\end{frame}

\begin{frame}{Why Google switched phones to ChaCha20}
  \begin{columns}[c,onlytextwidth]
    \column{0.3\textwidth}
      \centering
      \begin{tikzpicture}
        \draw[nInk, very thick, fill=white, rounded corners=6pt] (0,0) rectangle (1.8,3.2);
        \draw[nLine, fill=nFill] (0.15,2.55) rectangle (1.65,2.9);
        \node[font=\tiny, text=nInk] at (0.9,2.72) {\faLock\ https://};
        \foreach \y in {2.2,1.9,1.6,1.3} \draw[nLight, thick] (0.25,\y) -- (1.55,\y);
        \fill[nInk] (0.9,0.25) circle (0.1);
      \end{tikzpicture}
    \column{0.67\textwidth}
      \small \textbf{2014}: Chrome on Android starts to use ChaCha20 for \textbf{HTTPS} to Google.\\[6pt]
      Reason given by Google:\\[3pt]
      about 3 times faster than AES on devices without AES instructions, and RC4 was already broken.
  \end{columns}
  \vfill
  \keymsg{Without AES instructions, ChaCha20 is the fast choice.}
  % source: Google Security Blog, April 2014
  \say{Wpis na blogu Google Security z kwietnia 2014 roku; Google porównywał z AES w wersji używanej w HTTPS. HTTPS to zwykłe strony WWW z kłódką w pasku adresu.}
\end{frame}

% ======================================================================
\section[ChaCha20-Poly1305]{ChaCha20-Poly1305}
% ======================================================================

\begin{frame}{Recall: encryption does not protect integrity}
  \vfill
  \begin{center}
    \begin{tikzpicture}[>=Stealth]
      \node[rowlabel, font=\small] at (-0.5,0) {plaintext};
      \byterow{0}{0}{ptb}{P,a,y,{\textvisiblespace},1,0,0,{\textvisiblespace},E,U,R}{0.68}
      \node[rowlabel, font=\small] at (-0.5,-0.9) {ciphertext};
      \foreach \k in {0,...,10} \node[ctb] (c\k) at ({0.68*\k},-0.9) {};
      \draw[cellmark] (c4.south west) rectangle (c4.north east);
      \node[circle, fill=sPink, text=white, minimum size=0.75cm, font=\small] (x) at ({0.68*4},-2.0) {\faUserSecret};
      \draw[data, line width=2pt] (x) -- (c4.south);
      \node[font=\scriptsize, text=nInk, anchor=west] at ({0.68*4+0.5},-2.0) {the attacker XORs one byte with \texttt{0x08}};
      \node[rowlabel, font=\small] at (-0.5,-3.0) {decrypted};
      \byterow{0}{-3.0}{ptb}{P,a,y,{\textvisiblespace},9,0,0,{\textvisiblespace},E,U,R}{0.68}
    \end{tikzpicture}
  \end{center}
  \vfill
  \defbox{\textbf{Integrity}: the certainty that nobody changed the message.}
  \keymsg{Encryption hides the message, but does not stop changes.}
  \say{Przypomnienie: XOR szyfrogramu z wartością zmienia dokładnie ten sam bit tekstu jawnego. Odbiorca widzi poprawną wiadomość, klucz nie był potrzebny.}
\end{frame}

\samedot
\begin{frame}{A tag}
  \begin{center}\small The sender computes a short \textbf{tag} from the ciphertext and a secret key.\end{center}
  \vfill
  \begin{center}
    \begin{tikzpicture}[>=Stealth]
      \envelope{e}{0,0}{ciphertext}
      \node[circle, fill=nInk, text=white, font=\normalsize, minimum size=0.7cm] at (1.1,0.6) {\faStamp};
      \node[font=\scriptsize, text=nInk, anchor=north] at (0,-0.8) {+ tag (16 bytes)};
      \draw[data, line width=1.6pt] (1.5,0) -- (3.6,0) node[midway, above, font=\scriptsize] {sent};
      \node[actor=sTeal, font=\small\bfseries] (r) at (4.8,0) {\faServer\ receiver};
      \node[font=\scriptsize, text=nInk, anchor=west, align=left] at (6.0,0.45) {computes the tag again:};
      \node[font=\small, text=nInk, anchor=west] at (6.0,-0.05) {\ok\ same tag: accept and decrypt};
      \node[font=\small, text=nInk, anchor=west] at (6.0,-0.55) {\nok\ different tag: reject, decrypt nothing};
    \end{tikzpicture}
  \end{center}
  \vfill
  \begin{center}\small \textbf{Poly1305} is the function that computes such a tag.\end{center}
  \vfill
  \keymsg{A changed byte gives a different tag.}
  \say{Bez klucza napastnik nie potrafi policzyć nowego znacznika dla zmienionej wiadomości. Poly1305 to funkcja, która liczy taki znacznik.}
\end{frame}

\samedot
\begin{frame}{Authenticated encryption}
  \vfill
  \bigdef{\textbf{Authenticated encryption} (\textbf{AEAD}): encryption and a tag in one operation.\\
    Decryption first checks the tag, and refuses to decrypt if it is wrong.}
  \vfill
  \begin{center}
    \begin{tikzpicture}
      \node[hwmain, minimum width=3.2cm, minimum height=1.2cm, font=\small\bfseries, align=center] at (0,0) {ChaCha20\\[-1pt]{\scriptsize\normalfont for secrecy}};
      \node[font=\Large] at (2.3,0) {$+$};
      \node[hwmain, minimum width=3.2cm, minimum height=1.2cm, font=\small\bfseries, align=center] at (4.6,0) {Poly1305\\[-1pt]{\scriptsize\normalfont for integrity}};
      \node[font=\Large] at (6.9,0) {$=$};
      \node[hw, minimum height=1.2cm, font=\small\bfseries] at (9.3,0) {ChaCha20-Poly1305};
    \end{tikzpicture}
  \end{center}
  \vfill
  \say{AEAD to skrót od authenticated encryption with associated data. Associated data to jawne dane, na przykład nagłówek, które też są chronione znacznikiem; dziś z nich nie korzystamy.}
\end{frame}

\samedot
\begin{frame}{Where the Poly1305 key comes from}
  \vfill
  \begin{center}
    \begin{tikzpicture}[>=Stealth]
      \node[draw=AccentOrange, thick, fill=AccentOrange!15, minimum width=1.8cm, minimum height=0.7cm, font=\scriptsize] (b0) at (0,0) {64 bytes};
      \node[font=\scriptsize\bfseries, text=nInk, anchor=south] at (0,0.4) {block 0};
      \foreach \n in {1,2,3} {
        \node[ksb, minimum width=1.8cm, minimum height=0.7cm, font=\scriptsize] at ({2.0*\n},0) {64 bytes};
        \node[font=\scriptsize\bfseries, text=nInk, anchor=south] at ({2.0*\n},0.4) {block \n};
      }
      \node[font=\large, text=nInk] at (7.4,0) {\ldots};
      \draw[data] (b0.south) -- (0,-1.0) node[below, font=\scriptsize, text=nInk, align=center] {first 32 bytes:\\ a one-time key for Poly1305};
      \draw[decorate, decoration={brace, mirror, amplitude=5pt}, nInk] (1.1,-0.45) -- (7.0,-0.45)
        node[midway, below=6pt, font=\scriptsize, text=nInk] {keystream for the message};
    \end{tikzpicture}
  \end{center}
  \vfill
  \begin{center}\small A new nonce gives a new Poly1305 key. One key per message.\end{center}
  \vfill
  \keymsg{This is why ChaCha20-Poly1305 starts encrypting at counter 1.}
  \say{Poly1305 jest bezpieczny tylko wtedy, gdy jego klucz jest użyty raz. ChaCha20 daje nowy klucz za darmo z bloku 0. To wyjaśnia licznik 1 z zadania 1.}
\end{frame}

\begin{frame}{Task 9: change one byte}
  \taskhead{Task 9}{Message:\ \texttt{Pay 100 EUR to Anna.}\ \ Encrypted twice: only ChaCha20, and ChaCha20-Poly1305.}
  \vfill
  \begin{center}
    \begin{tikzpicture}
      \envelope{a}{0,0}{no tag}
      \node[font=\scriptsize, text=nInk, anchor=north] at (0,-0.75) {only ChaCha20};
      \envelope{b}{4.0,0}{with a tag}
      \node[font=\scriptsize, text=nInk, anchor=north] at (4.0,-0.75) {ChaCha20-Poly1305};
      \node[circle, fill=nInk, text=white, font=\normalsize, minimum size=0.7cm] at (5.1,0.6) {\faStamp};
    \end{tikzpicture}
  \end{center}
  \vfill
  \begin{center}\small Change \texttt{"1"} into \texttt{"9"} in both ciphertexts (XOR byte 4 with \texttt{0x08})\\ and try to decrypt them.\end{center}
  \vfill
  \say{Około 5 minut. aead.py jest gotowy i napisany w czystym Pythonie, zgodny z wektorami z RFC 8439.}
\end{frame}

\samedot
\begin{frame}[fragile]{Task 9: what to do}
  \begin{columns}[T,onlytextwidth]
    \column{0.62\textwidth}
\begin{lstlisting}[style=py, emph={chacha20,seal,unseal,print}, basicstyle=\color{VSText}\ttfamily\fontsize{6.0}{7.0}\selectfont]
# Without a tag: only ChaCha20, the data starts at counter 1
plain = chacha20(key, nonce, 1, message)
# With a tag: ChaCha20-Poly1305
sealed = seal(key, nonce, message)

# The attacker changes "1" (0x31) into "9" (0x39): XOR 0x08 at position 4.
# TODO 1: make changed_plain: a copy of plain with byte 4 XORed with 0x08
#         hint: b = bytearray(plain); b[4] ^= 0x08; changed_plain = bytes(b)
# TODO 2: make changed_sealed in the same way from sealed
# TODO 3: decrypt changed_plain with chacha20(key, nonce, 1, ...) and print it
# TODO 4: decrypt changed_sealed with unseal(key, nonce, ...) inside
#         try: ... except TagError as e: print("rejected:", e)
\end{lstlisting}
      {\scriptsize\textcolor{gray}{Shortened on the slide: the import, the key, the nonce and one print are in the file.}}
    \column{0.35\textwidth}
      \taskhead{Task 9}{Open \texttt{tamper.py} from \ghlink{Task9}.}
      \stepbox{Ready}{\texttt{aead.py}: \texttt{chacha20}, \texttt{seal}, \texttt{unseal}, \texttt{TagError}}
  \end{columns}
  \say{try i except łapią błąd: unseal rzuca TagError, gdy znacznik się nie zgadza.}
\end{frame}

\samedot
\begin{frame}{Solution: change one byte}
  \vfill
  \begin{center}
    % CHECKED: Task9/solution/tamper.py
    \small
    \renewcommand{\arraystretch}{1.8}
    \begin{tabular}{l l}
      without a tag: & \nok\ \texttt{Pay 900 EUR to Anna.} \\
      with a tag: & \ok\ \texttt{rejected: the tag does not match: the message was changed} \\
    \end{tabular}
  \end{center}
  \vfill
  \keymsg{With Poly1305, the change is caught. Without it, the change goes through.}
  \say{Kod sprawdzony. Odbiorca bez znacznika nie ma jak zauważyć zmiany: 900 to poprawny napis.}
\end{frame}

\begin{frame}{The nonce rule still holds}
  \begin{center}\small The same key and nonce used twice with ChaCha20-Poly1305:\end{center}
  \vfill
  \begin{columns}[T,onlytextwidth]
    \column{0.48\textwidth}
      \probbox{the same keystream}{$\rightarrow$ P1 XOR P2 leaks}
    \column{0.48\textwidth}
      \probbox{the same Poly1305 key}{$\rightarrow$ the attacker can make valid tags}
  \end{columns}
  \vfill
  \begin{center}\small A tag does not fix a repeated nonce. It makes the damage worse.\end{center}
  \vfill
  \keymsg{One key and nonce pair, one message. Always.}
  \say{Z dwóch wiadomości z tym samym kluczem Poly1305 napastnik potrafi wyliczyć ten klucz i podrabiać znaczniki. Szczegóły matematyczne pomijamy.}
\end{frame}

\samedot
\begin{frame}{Good to know: XChaCha20}
  \vfill
  \begin{center}
    \begin{tikzpicture}
      \foreach \k in {0,...,11} \node[draw=sAmber, line width=1pt, fill=white, minimum width=0.42cm, minimum height=0.45cm, inner sep=0pt] at ({0.45*\k},0) {};
      \node[font=\small, text=nInk, anchor=west] at (5.6,0) {96 bits: choosing it at random is risky after many messages};
      \foreach \k in {0,...,23} \node[draw=sAmber, line width=1pt, fill=white, minimum width=0.42cm, minimum height=0.45cm, inner sep=0pt] at ({0.45*\k},-0.9) {};
      \node[font=\small, text=nInk, anchor=west] at (11.0,-0.9) {192 bits};
    \end{tikzpicture}
  \end{center}
  \vfill
  \begin{center}\small \textbf{XChaCha20}: a nonce of 192 bits, so a random nonce is safe.\\
    Used for example in libsodium and in WireGuard's cookies.\end{center}
  \vfill
  \say{Można pominąć. Przy losowym nonce 96-bitowym ryzyko powtórki rośnie z liczbą wiadomości, jak w paradoksie urodzin; przy 192 bitach jest pomijalne.}
\end{frame}

\begin{frame}{Where you use it: the web}
  \vfill
  \begin{center}
    \small
    \renewcommand{\arraystretch}{1.9}
    \begin{tabular}{c l l}
      \textcolor{nInk}{\Large\faLock} & \textbf{TLS 1.3} (2018) & \texttt{TLS\_CHACHA20\_POLY1305\_SHA256}, one of three that all should support \\
      \textcolor{nInk}{\Large\faChrome} & \textbf{Chrome and Android} & since 2014, on devices without AES instructions \\
    \end{tabular}
  \end{center}
  \vfill
  \defbox{\textbf{TLS}: the protocol that encrypts HTTPS.}
  \say{TLS to protokół, który szyfruje HTTPS. RFC 8446 (w 2026 roku zastąpione przez RFC 9846) zaleca, żeby każdy program z TLS 1.3 obsługiwał ChaCha20-Poly1305, obok dwóch wersji AES z uwierzytelnieniem. Przeglądarka wybiera go zwykle tam, gdzie nie ma instrukcji AES.}
\end{frame}

\samedot
\begin{frame}{Where you use it: VPN and SSH}
  \vfill
  \begin{center}
    \small
    \renewcommand{\arraystretch}{1.9}
    \begin{tabular}{c l l}
      \textcolor{nInk}{\Large\faNetworkWired} & \textbf{WireGuard} & a VPN, uses only ChaCha20-Poly1305 for its data \\
      \textcolor{nInk}{\Large\faTerminal} & \textbf{OpenSSH 6.5} (2014) & \texttt{chacha20-poly1305@openssh.com}, today the default \\
    \end{tabular}
  \end{center}
  \vfill
  \defbox{\textbf{VPN}: an encrypted tunnel between a computer and a network.\quad \textbf{SSH}: an encrypted login to a remote computer.}
  \say{WireGuard celowo nie pozwala wybrać innego szyfru. ChaCha20-Poly1305 jest pierwszy na domyślnej liście szyfrów OpenSSH (man ssh\_config, sprawdzone).}
\end{frame}

\samedot
\begin{frame}{Where you use it: phones and the kernel}
  \vfill
  \begin{center}
    \small
    \renewcommand{\arraystretch}{1.9}
    \begin{tabular}{c l l}
      \textcolor{nInk}{\Large\faMobile} & \textbf{Adiantum} (2019) & disk encryption on cheap Android phones, with ChaCha12 \\
      \textcolor{nInk}{\Large\faLinux} & \textbf{Linux 4.8} (2016) & \texttt{/dev/urandom} produces random numbers with ChaCha20 \\
    \end{tabular}
  \end{center}
  \vfill
  \say{/dev/urandom to systemowe źródło liczb losowych w Linuksie. Generator na ChaCha20 zastąpił tam starszą konstrukcję, bo jest szybszy.}
\end{frame}

\begin{frame}{RC4 and ChaCha20 side by side}
  \vfill
  \begin{center}
    \small
    \renewcommand{\arraystretch}{1.6}
    \begin{tabular}{l l l}
      \toprule
      & \textbf{RC4} & \textbf{ChaCha20-Poly1305} \\
      \midrule
      nonce & \nok\ none & \ok\ 96 bits, built in \\
      random access & \nok\ no & \ok\ yes, by counter \\
      known bias & \nok\ yes & \ok\ none for 20 rounds \\
      integrity & \nok\ none & \ok\ Poly1305 tag \\
      status & \nok\ forbidden in TLS & \ok\ in TLS 1.3, SSH, WireGuard \\
      \bottomrule
    \end{tabular}
  \end{center}
  \vfill
  \keymsg{Every weakness of RC4 is fixed in ChaCha20-Poly1305.}
  \say{Każdy wiersz to jedna rzecz z dzisiejszych zajęć.}
\end{frame}

\begin{frame}{Summary}
  \vfill
  \begin{center}
    \renewcommand{\arraystretch}{1.6}
    \begin{tabular}{cl}
      \bigbadge{1} & Key, nonce and counter give any 64-byte block of the keystream directly. \\
      \bigbadge{2} & Inside are only three operations on 32-bit words: add, rotate, XOR. \\
      \bigbadge{3} & ChaCha20: a 4 by 4 state, 20 rounds of quarter rounds, add the initial state. \\
      \bigbadge{4} & Rounds spread every bit; the final addition stops the rounds being undone. \\
      \bigbadge{5} & Fast without special hardware, and with Poly1305 it also protects integrity. \\
    \end{tabular}
  \end{center}
  \vfill
  \say{Po jednym zdaniu na każdy cel wykładu.}
\end{frame}


\end{document}
